% Created 2020-11-23 Mon 17:49
% Intended LaTeX compiler: pdflatex
\documentclass[a4paper,cleardoubleempty,BCOR1cm]{scrbook}
\usepackage[utf8]{inputenc}
\usepackage[T1]{fontenc}
\usepackage{graphicx}
\usepackage{grffile}
\usepackage{longtable}
\usepackage{wrapfig}
\usepackage{rotating}
\usepackage[normalem]{ulem}
\usepackage{amsmath}
\usepackage{textcomp}
\usepackage{amssymb}
\usepackage{capt-of}
\usepackage{hyperref}
% use to waste space:
% \documentclass[12pt,a4paper]{article}

% if you have this style and like it.
%\documentclass{acmsiggraph}
%\documentclass[review]{acmsiggraph}      % review
%\documentclass[widereview]{acmsiggraph}  % wide-spaced review
%\documentclass[preprint]{acmsiggraph}    % preprint

% define a \comment{this is a comment which can have linebreaks in it}
\newcommand{\comment}[1]{}
% \newcommand{\todo}[1]{\marginpar{\bf{#1}}}
\newcommand{\todo}[1]{{\color{red}\bf{TODO: #1}}}

\usepackage{amsthm}
\usepackage{amsfonts}
\usepackage{mathptmx}
\usepackage{stmaryrd}
\usepackage{wasysym} %?
\usepackage{esvect}
\usepackage{graphicx}
\usepackage{color}
\definecolor{rot}{RGB}{165,30,55} %rote Farbe
\graphicspath{{./images/}}
\usepackage{parskip}
\usepackage{dsfont}
\usepackage{pxfonts}
\usepackage{tikz}

% comment these two lines out if you don't want minion/myriad fonts.
%\usepackage[minionint,mathlf]{MinionPro}
% \renewcommand{\sfdefault}{Myriad-LF}
%\usepackage{Myriad}
\usepackage{fontspec}
\setmonofont{DejaVu Sans Mono}

% no page number on float pages, fixes problems with overlarge diagrams.
\usepackage{fancyhdr}
\pagestyle{fancy}
%\lhead{}
%\chead{}
%\rhead{}
%\lfoot{}
\fancyhf{}
\fancyhead[EL]{\nouppercase{\leftmark}}
\fancyhead[OR]{\nouppercase{\rightmark}}
\cfoot{}
%\fancyfoot[EL]{\iffloatpage{}{\thepage}}
%\fancyfoot[OR]{\iffloatpage{}{\thepage}}
\fancyfoot[EL]{\thepage}
\fancyfoot[OR]{\thepage}
\renewcommand{\headrulewidth}{0pt}
\renewcommand{\footrulewidth}{0pt}

%\usepackage{natbib}		% textual referencing
%\usepackage[numbers,super]{natbib}	% nice superscripts
%\bibliographystyle{chicago}	% shitty
\bibliographystyle{alpha}	% abbr names and year in \cite
%\bibliographystyle{agsm}	% australian, need natbib
%\bibliographystyle{kluwer}	% need natbib
%\bibliographystyle{apalike}	% lengthly
%\bibliographystyle{abbrv}	% minimal?

% use for german line breaking:
%\usepackage[ngerman]{babel}
\usepackage[T1]{fontenc}

% avoid us-style text color destruction:
\frenchspacing
\usepackage{microtype}

% have a nice framebox with border directly around the image:
\fboxsep 0pt
\newcommand{\fimg}[2]{\fbox{\includegraphics[width=#1]{#2}}}

\newtheorem{theorem}{Theorem}
\newtheorem{definition}{Definition}

\usepackage{listings}
\lstset{basicstyle=\scriptsize, mathescape}
\usepackage{minted}
\usepackage[boxruled]{algorithm2e}
%\usepackage{hyperref}
\usepackage{url}
\usepackage{subfig}

\def\code#1{{\tt{#1}}}

\usepackage{bussproofs}
\usepackage{xcolor}
\usepackage{lscape}

\newcommand{\id}[1]{\text{id}_{#1}}
\newcommand{\rec}{\text{rec}}
\newcommand{\recT}[2]{\rec\;#1\text{ to } #2\text{ where}}
\newcommand{\corec}{\text{corec}}
\newcommand{\TyCtx}{\enskip\textbf{TyCtx}}
\newcommand{\ParCtx}{\enskip\textbf{ParCtx}}
\newcommand{\Ctx}{\enskip\textbf{Ctx}}
\newcommand{\D}{\AxiomC{$\mathcal{D}$}\noLine}
\newcommand{\Di}[1]{\AxiomC{$\mathcal{D}_{#1}$}\noLine}
\newcommand{\topI}[1]{\RightLabel{\textbf{($\top$-I)}}\UnaryInfC{#1}}
\newcommand{\TyVarI}[1]{\RightLabel{\textbf{TyVar-I}}\BinaryInfC{#1}}
\newcommand{\TyVarWeak}[1]{\RightLabel{\textbf{(TyVar-Weak)}}\BinaryInfC{#1}}
\newcommand{\TyWeak}[1]{\RightLabel{\textbf{(Ty-Weak)}}\BinaryInfC{#1}}
\newcommand{\TyInstLabel}{\RightLabel{\textbf{(Ty-Inst)}}}
\newcommand{\TyInst}[1]{\TyInstLabel\BinaryInfC{#1}}
\newcommand{\TyInstTrinary}[1]{\TyInstLabel\TrinaryInfC{#1}}
\newcommand{\ParamAbstr}[1]{\RightLabel{\textbf{(Param-Abstr)}}\UnaryInfC{#1}}
\newcommand{\FPTy}{\RightLabel{\textbf{(FP-Ty)}}}
\newcommand{\Inst}[1]{\RightLabel{\textbf{(Inst)}}\BinaryInfC{#1}}
\newcommand{\Conv}[1]{\RightLabel{\textbf{(Conv)}}\BinaryInfC{#1}}
\newcommand{\Proj}[1]{\RightLabel{\textbf{(Proj)}}\UnaryInfC{#1}}
\newcommand{\TermWeak}[1]{\RightLabel{\textbf{(Term-Weak)}}\BinaryInfC{#1}}
\newcommand{\IndILabel}{\RightLabel{\textbf{(Ind-I)}}}
\newcommand{\IndI}[1]{\IndILabel\UnaryInfC{#1}}
\newcommand{\IndIBinary}[1]{\IndILabel\BinaryInfC{#1}}
\newcommand{\CoindE}[1]{\RightLabel{\textbf{(Coind-E)}}\UnaryInfC{#1}}
\newcommand{\IndE}{\RightLabel{\textbf{(Ind-E)}}}
\newcommand{\CoindI}{\RightLabel{\textbf{(Coind-I)}}}
\newcommand{\rat}{\rightarrowtriangle}
\newcommand{\graybox}[1]{\colorbox{gray}{#1}}
\newcommand{\Unit}{\text{Unit}}
\newcommand{\Maybe}{\text{Maybe}}
\newcommand{\Just}{\text{Just}}
\newcommand{\Nothing}{\text{Nothing}}
\newcommand{\nothing}{\text{nothing}}
\newcommand{\Conat}{\text{Conat}}
\newcommand*\circled[1]{\tikz[baseline=(char.base)]{
    \node[shape=circle,draw,inner sep=1pt] (char) {#1};}}

\newenvironment{scprooftree}[1]%
 {\gdef\scalefactor{#1}\begin{center}\proofSkipAmount \leavevmode}%
 {\scalebox{\scalefactor}{\DisplayProof}\proofSkipAmount \end{center} }
\newenvironment{changemargin}[2]{%
\begin{list}{}{%
\setlength{\topsep}{0pt}%
\setlength{\leftmargin}{#1}%
\setlength{\rightmargin}{#2}%
\setlength{\listparindent}{\parindent}%
\setlength{\itemindent}{\parindent}%
\setlength{\parsep}{\parskip}%
}%
\item[]}{\end{list}}
\author{Florian Engel}
\date{\today}
\title{Implementation of Type Theory based on dependent Inductive and Coinductive Types}
\hypersetup{
 pdfauthor={Florian Engel},
 pdftitle={Implementation of Type Theory based on dependent Inductive and Coinductive Types},
 pdfkeywords={},
 pdfsubject={},
 pdfcreator={Emacs 26.1 (Org mode 9.1.9)}, 
 pdflang={English}}
\begin{document}

\maketitle
\begin{tabular}{lr}
% \includegraphics[width=0.5\linewidth]{logo_sw} % logo bw
 \includegraphics[width=0.5\linewidth]{UT_WBMW_Rot_4C} % logo red
 & \hspace{0.2\linewidth}
 \parbox{0.5\linewidth}{
   \large\bf\textsf{\color{rot}{Mathematisch-\\Naturwissenschaftliche\\Fakultät\\\\}}
   \hspace{-.144cm}\normalsize\textsf{\color{rot}{Programmiersprachen}}
   \vspace{0.6cm}
 }
\end{tabular}

\vspace*{10ex}
Masterarbeit

{\huge\bf\textsf{Implementation of Type Theory based on Dependent inductive and
    coinductive types}}

\vspace*{30ex}

Eberhard Karls Universität Tübingen\\
Mathematisch-Naturwissenschaftliche Fakultät\\
Wilhelm-Schickard-Institut für Informatik\\
Programmiersprachen\\
Florian Engel,~ \verb+florian.engel@student.uni-tuebingen.de+,~ 2020

\vspace*{5ex}

\begin{tabular}{@{}l@{\hspace{2em}}l}
  Bearbeitungszeitraum:& von-bis \vspace*{5ex} \\
  Betreuer/Gutachter:& Prof. Dr. Klaus Ostermann, Universität Tübingen\\
  Zweitgutachter:& Prof. Dr. Reinhard Kahle, Universität Tübingen
\end{tabular}

\thispagestyle{empty}
\newpage

\chapter*{Selbstst\"andigkeitserkl\"arung}
Hiermit versichere ich, dass ich die vorliegende Masterarbeit selbst\"andig und
nur mit den angegebenen Hilfsmitteln angefertigt habe und dass alle Stellen,
die dem Wortlaut oder dem Sinne nach anderen Werken entnommen sind,
durch Angaben von Quellen als Entlehnung kenntlich gemacht worden sind.
Diese Masterarbeit wurde in gleicher oder \"ahnlicher Form in keinem anderen
Studiengang als Pr\"ufungsleistung vorgelegt.

\vspace*{8ex}
\hrule
\vspace*{2ex}
Florian Engel (Matrikelnummer 3860700), \today



\chapter*{Abstract}
  Dependent types are a useful tool to restrict types even further than types of
  strongly typed languages like Haskell. This gives us further type safety. With
  them, we can also prove theorems. Coinductive types allow us to define types by
  their observations rather than by their constructors. This is useful for
  infinite types like streams. In many common dependently typed languages, like
  Coq and Agda, we can define inductive types which depend on values and
  coinductive types but not coinductive types, which depend on values.

In this work, we will first give a survey of coinductive types in these
languages and then implement the type theory from \cite{basold2016type}. This
type theory has both dependent inductive types and dependent coinductive
types. In this type theory the dependent function space becomes definable.
This leads to a more symmetrical approach of coinduction in dependently typed
languages.

\setcounter{tocdepth}{2}
\tableofcontents

\chapter{Introduction}
\label{sec:orgdf5c36b}
In functional programming, we have functions that get input and produce
output. These functions don't depend on multiple values i.e. if there is no IO
involved, they produce for the same input always the same output. For example,
if we call a function $\mathtt{or}$\;on the values $\mathtt{true}$\;and $\mathtt{false}$\;we always get
$\mathtt{true}$. This makes the code more predictable.

The $\mathtt{or}$\;function should only be working on Booleans. To call it on strings
$\mathtt{'foo'}$\;and $\mathtt{'bar'}$\;wouldn't make sense i.e. there is no defined output for
these inputs. To prevent calls like these, some functional programming
languages introduced types. Types contain only certain values. For example, the
type for truth values contains only the values for true and false. In Haskell
we can define it like the following:
\begin{minted}[fontsize=\scriptsize]{haskell}
data Bool = True | False
\end{minted}
This says we can construct values of type $\mathtt{Bool}$\;with the constructors $\mathtt{True}$
and $\mathtt{False}$. These types which have constructors are called inductive types. We can
then define $\mathtt{or}$\;like this:
\begin{minted}[fontsize=\scriptsize]{haskell}
or :: Bool -> Bool -> Bool
or True _    = True
or _    True = True
or _    _    = False
\end{minted}
Here, we just list equations that define what the output for a given input is.
For example, in the first equation, we say if the first value is constructed
with the constructor $\mathtt{True}$, we give back $\mathtt{True}$.  We don't care about the
second value, therefore we write $\mathtt{\_}$.  We are matching on the construction of
the input values.  Therefore, we call this method pattern matching.
If we call this function somewhere in the code on values that aren't of type
$\mathtt{Bool}$,  Haskell won't compile our code.  Instead, it gives back a type error.

If we now want to change $\mathtt{Bool}$\;to a three-valued logic, we have to add a
third constructor to $\mathtt{Bool}$. After that, we have to change every function which
pattern matches on $\mathtt{Bool}$. If there are a lot of those kinds of functions,
this would be a lot of repetitive work. If Haskell would have coinductive types,
this could be a lot less work. Coinductive types are types that are,
contrary to inductive types, defined over their destruction. So we could define
$\mathtt{Bool}$\;over its destructors. These would be or, and, etc.

Through this work, we will explain coinductive types at the examples of streams
and functions. They will be generalized to partial streams and the Pi type in
dependently typed languages. Streams are lists that are infinitely long. They
are useful for modeling many IO interactions. For example, a chat of a text
messenger might be infinitely long. We can never know if the chat is finished.
This is of course limited by the hardware, but we are interested in abstract
models. Functions are used everywhere in functional programming. In most of
these languages, they are first-class objects. But in languages with
coinductive types, we can define them. If we only have inductive and
coinductive types, we get a symmetrical language. This is useful because
then we can change an inductive type to a coinductive one and vice versa. It
is straight forward to add functions which destruct an inductive type by
pattern matching on the constructor. But it is hard to add a new constructor.
Then, we add this constructor to every pattern matching on that type.
For coinductive types it's the other way around. For more on this, see
\cite{binder2019decomposition}. In the implemented syntax we can define streams
like the following:
\begin{lstlisting}
codata Stream$\langle$A : Set$\rangle$ : Set where Hd : Stream (succ @ k)
  Hd : Stream $\rat$ A
  Tl : Stream $\rat$ Stream
\end{lstlisting}
And functions like follows:
\begin{lstlisting}
codata Fun$\langle$A : Set, B : Set$\rangle$ : Set where
  Inst : (x : A) $\rat$ Fun $\rat$ B
\end{lstlisting}
We can generalize streams to partial streams as the following:
\begin{lstlisting}
codata PStr$\langle$A : Set$\rangle$ : (n : Conat) $\rat$ Set where
  Hd : (k : Conat) $\rat$ PStr (succ @ k) $\rat$ A
  Tl : (k : Conat) $\rat$ PStr (succ @ k) $\rat$ PStr @ k
\end{lstlisting}
And functions to the Pi type.
\begin{lstlisting}
codata Pi$\langle$A : Set, B : (x : A) $\rat$ Set$\rangle$ : Set where
  Inst : (x : A) $\rat$ Pi $\rat$ B @ x
\end{lstlisting}

The rest of this thesis is structured as follows:

\begin{itemize}
\item Chapter \ref{sec:org5d34b3c} shows how coinductive types can be defined.
Here, we will define the stream and function type, as well as some
functions on the stream.

\item We will see in Chapter \ref{sec:orgb5a19f8} how
coinductive types are defined in the dependently typed languages Coq and Agda.
We will see that we can define them positive or negative. We will show why
defining them positive leads to problems.

\item In Chapter \ref{sec:org33f714d} we
see how they are defined by \cite{basold2016type}. With this theory we can then
define coinductive types which depend on values.  But we can not define types
that depend on types.

\item We will then in Chapter \ref{sec:org8542692} explain how this theory is
implemented. Therefore, we need to rewrite the typing rules. It will also be
possible to define type schemata.

\item At last, we look at the examples from this paper \cite{basold2016type} in the implemented syntax.
Here, we will see the reduction steps for recursion and corecursion. We will
conclude this section with the example of partial streams, which is a
coinductive type that depends on a value.
\end{itemize}

\chapter{Coinductive Types}
\label{sec:org5d34b3c}
Inductive types are defined via their constructors.  Coinductive types on
the other hand are defined via their destructors.  In the paper \cite{abel2013copatterns}
functions, which have coinductive types as their output, are implemented via
copattern matching.  In that paper streams are defined like the following:

\begin{lstlisting}
record Stream A = { head : A,
                    tail : Stream A }
\end{lstlisting}

The $\mathtt{A}$\;in the definition should be a concrete type. The type system in the
paper doesn't have dependent types. What differentiates this from regular record
types (for example in Haskell) is the recursive field tail. So they call it a
recursive record. In a strict language without coinductive types we could never
instantiate such a type because to do this we already need something of type
$\mathtt{Stream\;A}$\;to fill in the field $\mathtt{tail}$. The paper defines
copattern matching to remedy this. With the help of copattern matching, we can define functions
that output expressions of type $\mathtt{Stream\;A}$. As an example, we look at the
definition of repeat. This function takes in a value of type $\mathtt{Nat}$\;and
generates a stream that just infinitely repeats it.

\begin{lstlisting}
repeat : Nat $\rat$ Stream Nat
head (repeat x) = x
tail (repeat x) = repeat x
\end{lstlisting}

As we can see copattern matching works via observations i.e. we define what
should be the output of the fields applied to the result of the function. These
fields are also called observers because we observe parts of the type.
Because inhabitants of $\mathtt{Stream}$\;are infinitely long we can't print out a
stream. Because of this we also consider each expression which has a type,
which is coinductive, as a value. To get a subpart of this value we use
observers. For example, we can look at the third value of $\mathtt{repeat\;2}$\;via \texttt{head
\;\;(tail\;(tail\;(repeat\;2)))}\;which should evaluate to 2. We can also implement a
function that looks at the nth. value. Here it is:

\begin{lstlisting}
nth : Nat $\rat$ Stream A $\rat$ A
nth 0     x = head x
nth (S n) x = nth n (tail x)
\end{lstlisting}

As you can see we use ordinary pattern matching on the left-hand side and
observers on the right-hand side. $\mathtt{nth\;3\;(repeat\;2)}$\;will output $\mathtt{2}$\;as expected.
Functions can also be defined via a recursive record.  It is defined as the
following:

\begin{lstlisting}
record A $\rat$ B = { apply : A $\leadsto$ B }
\end{lstlisting}

Here, we differentiate between our defined function $\mathtt{A\;\rat\;B}$\;and $\mathtt{\leadsto}$\;in the
destructor. Constructor applications or, as is the case here, destructor
applications are not the same as function applications, like in Haskell. In the
paper $\mathtt{f\;x}$\;means $\mathtt{apply\;f\;x}$. We will also use this convention in the
following. In fact, we already used it in the definitions of the functions
$\mathtt{repeat}$\;and $\mathtt{nth}$. $\mathtt{nth\;0\;x\;=\;head\;x}$\;is just a nested copattern. We can also write it
with $\mathtt{apply}$\;like so: $\mathtt{apply\;(apply\;nth\;0)\;x\;=\;head\;x}$. Here, we use currying.
So the first apply is the sole observer of type $\mathtt{Stream\;A\;\rat\;A}$\;and the second
of type $\mathtt{Nat\;\rat\;(Stream\;A\;\rat\;A)}$.

\chapter{Coinductive Types in dependent languages}
\label{sec:orgb5a19f8}
In this section, we will look at how coinductive types are implemented in
dependently typed languages. In dependently typed languages types can depend on
values. The classical example of such a type is the type for vectors. Vectors are like
lists, except their length is contained in their type. For example, a vector of
natural numbers of length 2 has type $\mathtt{Vec\;Nat\;2}$. This type depends on two
things. Namely the type $\mathtt{Nat}$\;and the value $\mathtt{2}$, which is itself of type $\mathtt{Nat}$.
We can define vectors in Coq as follows:
\begin{minted}[fontsize=\scriptsize]{coq}
Inductive Vec (A : Set) : nat -> Set :=
  | Nil : Vec A 0
  | Cons : forall {k : nat}, A -> Vec A k -> Vec A (S k).
\end{minted}
Contrary to a list the type constructor $\mathtt{Vec}$\;has a second argument $\mathtt{nat}$.
This is the already mentioned length of the vector. A Vector has two
constructors. One for an empty vector called $\mathtt{Nil}$\;and one to append an
element at the front of a vector called $\mathtt{Cons}$. $\mathtt{Nil}$\;just returns a vector
of length $\mathtt{0}$. And $\mathtt{Cons}$\;gets an $\mathtt{A}$\;and a vector of length $\mathtt{k}$. It returns a
vector of length $\mathtt{S\;k}$\;($\mathtt{S}$\;is just the successor of k). This type can also be
defined in agda as follows:
\begin{minted}[fontsize=\scriptsize]{agda}
data Vec (A : Set) : ℕ → Set where
  Nil : Vec A 0
  Cons : {k : ℕ} → A → Vec A k → Vec A (suc k)
\end{minted}
One advantage of vectors in comparison to lists is that we can define a total function
(a function which is defined for every input) that takes the head of a
vector. This function can't be total for lists, because we cannot know if the
input list is empty. An empty list has no head. For vectors, we can enforce this
in Coq like follows:
\begin{minted}[fontsize=\scriptsize]{coq}
Definition hd {A : Set} {k : nat} (v : Vec A (S k)) : A :=
  match v with
  | Cons _ x _ => x
  end.
\end{minted}
We just pattern match on $\mathtt{v}$.  The only pattern is for the $\mathtt{Cons}$\;constructor.  The $\mathtt{Nil}$\;constructor
is a vector of length 0.  But $\mathtt{v}$\;has type $\mathtt{Vec\;A\;(S\;k)}$.  So it can't be a vector of length 0.
In Agda the function looks like follows:
\begin{minted}[fontsize=\scriptsize]{agda}
hd : {A : Set} {k : ℕ} → Vec A (suc k) → A
hd (cons x _) = x
\end{minted}
That types can depent on terms makes it necessary to ensure that function
\linebreak terminate. Otherwise, type checking wouldn't be decidable. If we have a
function \linebreak $\mathtt{f\;:\;Nat\;\rat\;Nat}$\;and we want to check a value $\mathtt{a}$\;against a type
$\mathtt{Vec\;(f\;1)}$\;we have to know what $\mathtt{f\;1}$\;evaluates to. So $\mathtt{f}$\;has to terminate.
We check termination in Coq via a structural decreasing argument. An argument
is structural decreasing if it is structural smaller in a recursive call.
Structural smaller means it is a recursive occurrence in a constructor. As an
example, we look at the definition of the natural numbers and the function for addition
on them. We define the natural numbers in Coq like follows:
\begin{minted}[fontsize=\scriptsize]{coq}
Inductive nat : Set :=
| O : nat
| S : nat -> nat.
\end{minted}
$\mathtt{O}$\;is the constructor for 0 and $\mathtt{S}$\;is the successor of its argument. Here,
the recursive argument to $\mathtt{S}$\;is structurally smaller than $\mathtt{S}$\;applied to it i.e.
$\mathtt{n}$\;is structurally smaller than $\mathtt{S\;n}$. Then, we can define addition like follows:
\begin{minted}[fontsize=\scriptsize]{coq}
Fixpoint add (n m : nat) : nat :=
match n with
| O => m
| S p => S (add p m)
end.
\end{minted}
In the recursive call, the first argument is structural decreasing. The
expression $\mathtt{p}$\;is smaller than the expression $\mathtt{s\;p}$. So Coq accepts this
definition. The classical example of a function where an argument is
decreasing but not structural decreasing is Quicksort. A naive implementation
would be the following:
\begin{minted}[fontsize=\scriptsize]{coq}
Fixpoint quicksort (l : list nat) : list nat :=
match l with
| nil => nil
| cons x xs => match split x xs with
              | (lower, upper) => app (quicksort lower) (cons x (quicksort upper))
              end
end.
\end{minted}
Here, $\mathtt{split}$\;is just a function that gets a number and a list of numbers. It
gives back a pair of two lists where the elements of the left list are all
elements of the input list which are smaller than the input number and the
right these which are bigger. It is clear that these lists can't be longer
than the input list. So $\mathtt{lower}$\;and $\mathtt{upper}$\;can't be longer than $\mathtt{xs}$. Here
$\mathtt{xs}$\;is structurally smaller than the input $\mathtt{cons\;x\;xs}$. So $\mathtt{lower}$\;and $\mathtt{upper}$
are smaller than the input. Therefore, we know that $\mathtt{quicksort}$\;is terminating.
But Coq won't accept our code, because no argument is structural decreasing.

For coinductive types termination means that functions that produce them
should be productive. If a function is productive it produces in each step a
new part of the infinitely large coinductive type.

In Section \ref{sec:org5d9a58e} we will look at the implementation in Coq.
There are two ways to define them. The older way uses positive coinductive
types. This is known to violate subject reduction. Therefore, it is highly
discouraged to use them. To fix this the new way uses negative coinductive
types. In Section \ref{sec:org9447cd3} we look at the implementation in
Agda. Agda also has these two ways of defining such types. One special thing
about it, is that it implements copattern matching. To help Agda with
termination checking we can use sized types.  We will explain them in Section
\ref{sec:org5ddc9c0}.
\section{Coinductive Types in Coq}
\label{sec:org5d9a58e}
There are two approaches to define coinductive types in Coq. The older one is
described in \ref{sec:org1db9e9e}. It works over constructors. Therefore
they are called positive coinductive types. The newer and recommended one is
described in Section \ref{sec:orgf9d728b}. They are defined over
primitive records (a relatively new feature of Coq). Therefore, they are
called negative coinductive Types.

\subsection{Positive Coinductive Types}
\label{sec:org1db9e9e}
Positive coinductive types are defined over constructors in Coq.  The keyword
\linebreak $\mathtt{CoInductive}$\;is used to indicate that we about to define a coinductive type.
This is the only syntactical difference from the definition of inductive
types. For example, streams are defined like the following:

\begin{minted}[fontsize=\scriptsize]{coq}
CoInductive Stream (A : Set) : Set :=
  Cons : A -> Stream A -> Stream A.
\end{minted}

If this was an inductive type we couldn't generate a value of this type. To
generate values of coinductive types Coq uses guarded recursion. This checks
if the recursive call to the function occurs as an argument to a coinductive
constructor. In addition to the guard condition, the constructor can only be
nested in other constructors, fun or match expressions. With all of this in
mind we can define $\mathtt{repeat}$\;like the following:

\begin{minted}[fontsize=\scriptsize]{coq}
CoFixpoint repeat (A : Set) (x : A) : Stream A := Cons A x (repeat A x).
\end{minted}

Then, we can produce the constant zero stream with $\mathtt{repeat\;nat\;0}$. If we used
a normal Coq function i.e. write $\mathtt{Fixpoint}$\;instead of $\mathtt{CoFixpoint}$\;Coq
wouldn't accept our code. It rejects it because there is no argument which is
structural decreasing. $\mathtt{x}$\;stays always the same. Functions defined with
$\mathtt{CoFixpoint}$\;on the other hand only check the previously mentioned
conditions. It sees the recursive call $\mathtt{repeat\;A\;x}$\;occurs as an argument to
constructor $\mathtt{Cons}$\;of the coinductive type $\mathtt{Stream}$. This constructor is also
not nested. So our definition is accepted.

We can use the normal pattern matching of Coq to destruct a coinductive type.
We define $\mathtt{nth}$\;like the following:

\begin{minted}[fontsize=\scriptsize]{coq}
Fixpoint nth (A : Set) (n : nat) (s : Stream A) {struct n} : A :=
  match s with
    Cons _ a s' =>
    match n with 0 => a | S p => nth A p s' end
  end.
\end{minted}

The guard condition is necessary to ensure every expression is terminating.
If we didn't have the guard condition we could define the following:

\begin{minted}[fontsize=\scriptsize]{coq}
CoFixpoint loop (A : Set) : Stream A = loop A.
\end{minted}

Here, the recursive call doesn't occur in a constructor. So the guard
condition is violated. With this definition the expression $\mathtt{nth\;0\;loop}$
wouldn't terminate. The function $\mathtt{nth}$\;would try to pattern match on $\mathtt{loop}$. But to
succeed in that $\mathtt{loop}$\;has to unfold to something of the form $\mathtt{Cons\;a\;?}$
which it never does. So $\mathtt{nth\;0\;loop}$\;will never evaluate to a value. This
would lead to undecidable type checking.

We illustrate the purpose of the other conditions on an example taken from
\cite{chlipala2013certified}.  First, we implement the function $\mathtt{tl}$\;like so:

\begin{minted}[fontsize=\scriptsize]{coq}
Definition tl A (s : Stream A) : Stream A :=
  match s with
  | Cons _ _ s' => s'
  end.
\end{minted}

This is just one normal pattern match on $\mathtt{Stream}$.  If we didn't have the
other condition we could define the following:

\begin{minted}[fontsize=\scriptsize]{coq}
CoFixpoint bad : Stream nat := tl nat (Cons nat 0 bad).
\end{minted}

This doesn't violate the guard condition.  The recursive call $\mathtt{bad}$\;is an
argument to the constructor $\mathtt{Cons}$.  But the constructor is nested in a
function.  If we would allow this, $\mathtt{nth\;0\;bad}$\;would loop forever.  To
understand why we first unfold $\mathtt{tl}$\;in $\mathtt{bad}$.  So we get:

\begin{minted}[fontsize=\scriptsize]{coq}
nth 0 (cofix bad : Stream nat :=
         match (Cons 0 bad) with
         | Cons _ s' => s'
         end)
\end{minted}

We can now simplify this to just:

\begin{minted}[fontsize=\scriptsize]{coq}
nth 0 (cofix bad : Stream nat := bad)
\end{minted}

After that $\mathtt{bad}$\;isn't any more an argument to a constructor.  Here, we can also
see easily that the expression $\mathtt{cofix\;bad\;:\;Stream\;nat\;:=\;bad}$\;loops forever.
So we never get the value at position $\mathtt{0}$.

An important property of typed languages is subject reduction. Subject
reduction says if we evaluate an expression \(e_1\) of type \(t\) to an expression
\(e_2\), \(e_2\) should also be of type \(t\). With positive coinductive types subject
reduction is no longer valid. We illustrate this by Oury's counterexample
\cite{oury2008}. First, we define the codata type $\mathtt{U}$\;as follows:

\begin{minted}[fontsize=\scriptsize]{coq}
CoInductive U : Set := In : U -> U.
\end{minted}

We can now define a value of u with the following $\mathtt{Cofixpoint}$\;like so:

\begin{minted}[fontsize=\scriptsize]{coq}
CoFixpoint u : U := In u.
\end{minted}

This generates an infinite succession of $\mathtt{In}$.  We use the function $\mathtt{force}$
to force $\mathtt{U}$\;to evaluate one step i.e. $\mathtt{x}$\;becomes $\mathtt{In\;y}$.

\begin{minted}[fontsize=\scriptsize]{coq}
Definition force (x: U) : U :=
  match x with
    In y => In y
  end.
\end{minted}

The same trick will be used to define $\mathtt{eq}$\;which states that $\mathtt{x}$\;is
definitional equal to $\mathtt{force\;x}$.

\begin{minted}[fontsize=\scriptsize]{coq}
Definition eq (x : U) : x = force x :=
  match x with
    In y => eq_refl
  end.
\end{minted}

This first matches on $\mathtt{x}$\;to force it, to reduce to $\mathtt{In\;y}$. Then, the new goal
becomes $\mathtt{In\;y\;=\;force\;(In\;y)}$. $\mathtt{force\;(In\;y)}$\;evaluates to just $\mathtt{In\;y}$, as it
is just pattern matching on $\mathtt{In\;y}$. So the final goal is $\mathtt{In\;y\;=\;In\;y}$\;which
can be shown by $\mathtt{eq\_refl}$. The expression $\mathtt{eq\_refl}$\;is a constructor for $\mathtt{=}$
where both sides of $\mathtt{=}$\;are exactly the same. If we now instantiate $\mathtt{eq}$\;with
$\mathtt{u}$\;we become $\mathtt{eq\;u}$.

\begin{minted}[fontsize=\scriptsize]{coq}
Definition eq_u : u = In u := eq u
\end{minted}

But $\mathtt{u}$\;is not definitional equal to $\mathtt{In\;u}$.  As mentioned above expressions
with a coinductive type are always values to prevent infinite evaluation.
So $\mathtt{In\;u}$\;is a value and $\mathtt{u}$\;is also a value.  But values are only
definitional equal if they are exactly the same.  The next section will
solve this problem through negative coinductive types.

\subsection{Negative Coinductive Types}
\label{sec:orgf9d728b}
In Coq 8.5. primitive records were introduced. With this, it is now possible
to define types over their destructors. So we can have negative, especially
negative coinductive, types in Coq. With primitive records we can define
streams like the following:

\begin{minted}[fontsize=\scriptsize]{coq}
CoInductive Stream (A : Set) : Set :=
  Seq { hd : A; tl : Stream A }.
\end{minted}

Now we can define $\mathtt{repeat}$\;over the fields of $\mathtt{Stream}$.

\begin{minted}[fontsize=\scriptsize]{coq}
CoFixpoint repeat (A : Set) (x : A) : Stream A :=
  {| hd := x; tl := repeat A x|}.
\end{minted}

To define $\mathtt{repeat}$\;we must define what is the head of the constructed stream
and its tail.  The guard condition says now that corecursive
occurrences must be guarded by a record field.  We can see that the
corecursive call $\mathtt{repeat}$\;is a direct argument to the field $\mathtt{tl}$\;of the
corecursive type $\mathtt{Stream\;A}$.  This means Coq accepts the above definition.
If we want to access parts of a stream we use the destructors $\mathtt{hd}$\;and
$\mathtt{tl}$.  With them, we can define $\mathtt{nth}$\;again for the negative stream.

\begin{minted}[fontsize=\scriptsize]{coq}
Fixpoint nth (A : Set) (n : nat) (s : Stream A) : list A :=
  match n with
  | 0 => s.(hd A)
  | S n' => nth A n' s.(tl A)
  end.
\end{minted}

With negative coinductive types, we can't form the above-mentioned
counterexample to subject reduction anymore, because we can't pattern match
on negative types. Oury's example becomes.

\begin{minted}[fontsize=\scriptsize]{coq}
CoInductive U := { out : U }.
\end{minted}

$\mathtt{U}$\;is now defined over its destructor $\mathtt{out}$, instead of its constructor
 $\mathtt{in}$. Then, $\mathtt{in}$\;\linebreak becomes just a function. In fact, it's just a
 definition because we don't recurse or corecurse on it.

\begin{minted}[fontsize=\scriptsize]{coq}
Definition In (y : U) : U := {| out := y |}.
\end{minted}

We define it over the only field $\mathtt{out}$.  When we put a $\mathtt{y}$\;in then we get
the same $\mathtt{y}$\;out.  We can also again define $\mathtt{u}$.

\begin{minted}[fontsize=\scriptsize]{coq}
CoFixpoint u : U := {| out := u |}.
\end{minted}

With coinductive types, it is now possible to define the pi type (the dependent
function type).

\begin{minted}[fontsize=\scriptsize]{coq}
CoInductive Pi (A : Set) (B : A -> Set) := { Apply (x : A) : B x }.
\end{minted}

The pi type is defined over its destructor $\mathtt{Apply}$.  If we evaluate $\mathtt{Apply}$
on a value of $\mathtt{Pi}$\;(which is a function) and an argument, we get the result
i.e. we apply the value to the function.  It looks like the pi type becomes definable
in Coq.  But we are cheating.  The type of $\mathtt{Apply}$\;is already a pi type because
we identify constructors and destructors with functions.  We will see that the theory of
the paper avoids this identification. To define a function we use
$\mathtt{CoFixpoint}$.  As a simple nonrecursive, nondependent example we use the
function $\mathtt{plus2}$.

\begin{minted}[fontsize=\scriptsize]{coq}
CoFixpoint plus2 : Pi nat (fun _ => nat) :=
  {| Apply x  := S (S x) |}.
\end{minted}

If we apply (i.e. call the destructor $\mathtt{Apply}$) a $\mathtt{x}$\;to $\mathtt{plus2}$\;it gives back
$\mathtt{S\;(S\;x)}$.  Which is twice the successor on $\mathtt{x}$.  So we add 2 to $\mathtt{x}$.  We
use $\mathtt{\_}$\;here because $\mathtt{plus2}$\;is not a dependent function i.e. the result
type $\mathtt{nat}$\;doesn't depend on the input value.  To define functions with more
than one argument we just use currying i.e. we use the type $\mathtt{Pi}$\;as the
second argument to $\mathtt{Pi}$. For example, a 2-ary non-dependent function from $\mathtt{A}$
and $\mathtt{B}$\;to $\mathtt{C}$\;would have type $\mathtt{Pi\;A\;(fun\;\_\;=>\;Pi\;B\;(fun\;\_\;=>\;C))}$.  It
would be fortunate if we could define $\mathtt{plus}$\;like the following:

\begin{minted}[fontsize=\scriptsize]{coq}
CoFixpoint plus : Pi nat (fun _ => Pi nat (fun _ => nat)) :=
  {| Apply := fun (n : nat)  =>
       match n with
       | O => {| Apply (m : nat) := m |}
       | S n' => {| Apply m := S (Apply _ _ (Apply _ _  plus n') m) |}
       end
  |}.
\end{minted}

But Coq doesn't accept this definition. The guard condition is violated. The
expression $\mathtt{plus\;n'}$\;is not a direct argument of the field $\mathtt{Apply}$. The
definition should terminate because we are decreasing $\mathtt{n}$\;and the case for
$\mathtt{0}$\;is accepted. In the case of $\mathtt{0}$, there is no recursive call.

We can also define a dependent function.  We define append2Units like
follows
\begin{minted}[fontsize=\scriptsize]{coq}
CoFixpoint append2Units : Pi nat
                             (fun n => Pi (Vec unit n)
                                       (fun _ => Vec unit (S (S n)))) :=
  {| Apply n := {| Apply v := Cons _ tt (Cons _ tt v) |} |}.
\end{minted}
This just appends 2 units at a vector of length $\mathtt{n}$.  Here, the second
argument and the result depend on the first argument i.e. the first argument
is the length of the input vector and the output vector is this length plus
two.

\section{Coinductive Types in Agda}
\label{sec:org9447cd3}
In Agda coinductive types were first also introduced as positive types. In
Section \ref{sec:org6e1a082} we will look at them in
detail. In Section \ref{sec:orgda75ba9} we describe the correct
way to implement coinductive types in Agda. There are functions which
terminate but are rejected by the type checker.   To allow more functions we can use a unique feature of Agda, sized
types. They are described in Section \ref{sec:org5ddc9c0}.

\subsection{Positive Coinductive Types in Agda}
\label{sec:org6e1a082}
Agda doesn't have a special keyword to define coinductive types like Coq.  It
uses the symbol $\mathtt{\infty}$\;to mark arguments to constructors as coinductive.
This symbol says that the computation of arguments of this type are suspended.
$\mathtt{\infty}$\;is just a type constructor.  So Agda ensures productivity over type
checking. We define streams like so.

\begin{minted}[fontsize=\scriptsize]{agda}
data Stream (A : Set) : Set where
  cons : A → ∞ (Stream A) → Stream A
\end{minted}

Here, the second argument to $\mathtt{cons}$\;is marked with $\mathtt{\infty}$. This is the tail of
the stream. Because it is infinitely long (we don't have a constructor of an
empty stream) we can't compute it completely, so we suspend the computation.
We can delay a computation with the constructor $\mathtt{\sharp}$\;and force it with
the function $\mathtt{\flat}$. Their types are given below.

\begin{minted}[fontsize=\scriptsize]{agda}
♯_ : ∀ {a} {A : Set a} → A → ∞ A
♭  : ∀ {a} {A : Set a} → ∞ A → A
\end{minted}

We can now again define our usual functions.  We begin with $\mathtt{repeat}$.

\begin{minted}[fontsize=\scriptsize]{agda}
repeat : {A : Set} → A → Stream A
repeat x = cons x (♯ (repeat x))
\end{minted}

We first apply $\mathtt{cons}$\;to $\mathtt{x}$. So the head of the stream is $\mathtt{x}$. We then apply
it to the corecursive call $\mathtt{repeat}$. So the tail will be a repetition of
$\mathtt{xs}$. We have to call the $\mathtt{repeat}$\;with $\mathtt{\sharp}$\;to suspend the computation.
Otherwise, the code doesn't type check. If we would write this function
without $\mathtt{\sharp}$\;on a stream which has no $\mathtt{\infty}$\;on the second argument of
$\mathtt{cons}$, the function would run forever. In fact, the termination checker won't
allow us to write such a function. We can also write $\mathtt{nth}$\;again, which
consumes a stream.

\begin{minted}[fontsize=\scriptsize]{agda}
nth : {A : Set} → ℕ → Stream A → A
nth 0       (cons x _)  = x
nth (suc n) (cons _ xs) = nth n (♭ xs)
\end{minted}

Here, we have to use $\mathtt{\flat}$\;on the right-hand side of the second case, to
force the computation of the tail of the input stream.  We have to do that
because $\mathtt{nth}$\;wants a stream.  It doesn't want a suspended stream.
Productivity on coinductive types like stream is checked by only allowing non
decreasing recursive calls behind the $\mathtt{\sharp}$\;constructor.

\subsection{Negative Coinductive Types in Agda}
\label{sec:orgda75ba9}
In Agda we can also define negative coinductive types.  This is the
recommended way.  Agda implements the previously mentioned copattern matching.
We can define a record with the keyword $\mathtt{record}$.  We use the keyword $\mathtt{coinductive}$
to make it possible to define recursive fields.  Stream is defined as the
following:

\begin{minted}[fontsize=\scriptsize]{agda}
record Stream (A : Set) : Set where
  coinductive
  field
    hd : A
    tl : Stream A
\end{minted}

A Stream has 2 fields. The field $\mathtt{hd}$\;is the head of the stream. It has type
$\mathtt{A}$. The field $\mathtt{tl}$\;is the tail of the stream. It is another stream, so it
has type $\mathtt{Stream\;A}$. $\mathtt{tl}$\;is a recursive field. So Agda wouldn't accept the
definition without $\mathtt{coinductive}$. A stream can never be empty. Every stream
has a head (a field $\mathtt{hd}$) and an empty stream wouldn't have a head. So the
tail of a stream can never be empty. Therefore, every stream is infinitely
long. We can now define $\mathtt{repeat}$\;with copattern matching.

\begin{minted}[fontsize=\scriptsize]{agda}
repeat : ∀ {A : Set} → A → Stream A
hd (repeat x) = x
tl (repeat x) = repeat x
\end{minted}

We have to copattern match on every field of $\mathtt{Stream}$, namely $\mathtt{hd}$\;and $\mathtt{tl}$.
Because Agda is total it won't accept non-exhaustive (co)pattern matches
like Haskell.  First, we define what the head of $\mathtt{repeat\;x}$\;is.  We just
repeat $\mathtt{x}$\;infinitely often.  So every element of the steam is $\mathtt{x}$, including
the head.  Therefore, we just write $\mathtt{x}$.  In the second and last copattern we
define what the tail of the stream is.  The tail is just $\mathtt{repeat\;x}$.
Infinitely often repeated $\mathtt{x}$\;is the same as x and then infinitely repeated
$\mathtt{x}$.  We can use normal pattern matching and the destructors for functions
that consume streams.  We define $\mathtt{nth}$\;like the following:

\begin{minted}[fontsize=\scriptsize]{agda}
nth : ∀ {A : Set} → ℕ → Stream A → A
nth zero s = hd s
nth (suc n) s = nth n (tl s)
\end{minted}

Here, we just pattern match on the first argument (excluding the implicit
argument of the type).  If it is zero the result is just the head of the
stream.  If it is \(n+1\) the result is the recursive call of $\mathtt{nth}$\;on $\mathtt{n}$\;and
$\mathtt{tl\;s}$.  Agda accepts this code because it is structural decreasing on the
first (or second if we count the implicit) argument.

We can also define the pi type.  We use $\mathtt{\_\$\_}$\;as the apply operator.  This
operator is taken from Haskell.

\begin{minted}[fontsize=\scriptsize]{agda}
record Pi (A : Set) (B : A → Set) : Set where
  field _$_ : (x : A) → B x
  infixl 20 _$_
open Pi
\end{minted}

Like in Coq we are using the first-class pi type to define the pi type. Agda
doesn't define the first-class pi type lkie that.  We
can also define a function which adds 2 to a number $\mathtt{plus2}$\;in Agda.

\begin{minted}[fontsize=\scriptsize]{agda}
plus2 : ℕ →' ℕ
plus2 $ x = suc (suc x)
\end{minted}

We just use copattern matching to define it. If we apply a $\mathtt{x}$\;to $\mathtt{plus2}$\;we
get \linebreak $\mathtt{suc\;(suc\;x)}$. $\mathtt{\rightarrow'}$\;is just the non-dependent function it is defined
using our pi type. Here it is:

\begin{minted}[fontsize=\scriptsize]{agda}
_→'_ : Set → Set → Set
A →' B = Pi A (λ _ → B)
infixr 20 _→'_
\end{minted}

In Agda it becomes possible to define plus. We just use nested copattern
matching.

\begin{minted}[fontsize=\scriptsize]{agda}
plus : ℕ →' ℕ →' ℕ
plus $ 0       $ m = m
plus $ (suc n) $ m = suc (plus $ n $ m)
\end{minted}

If we change $\mathtt{\rightarrow'}$\;to $\mathtt{\rightarrow}$\;and remove $\mathtt{\$}$\;we get the
standard definition for plus in Agda. We can also define a dependent function
$\mathtt{repeatUnit}$\;like follow:
\begin{minted}[fontsize=\scriptsize]{agda}
repeatUnit : Pi ℕ (λ n → Vec ⊤ n)
repeatUnit $ 0     = nil
repeatUnit $ suc n = tt :: (repeatUnit $ n)
\end{minted}
This function gives back a vector with the length of the input, where every element
is unit.

\subsection{Termination Checking with Sized Types}
\label{sec:org5ddc9c0}
They are many functions which are total but are not accepted by Agda's
termination checker. In fact, in any total language, there have to be such
functions. We can show that by trying to list all total functions. The
following table lists functions per row. The columns say what the output of
the functions for the given input is.
\begin{center}
\begin{tabular}{lrrrrl}
 & 1 & 2 & 3 & 4 & \(\dots\)\\
\hline
\(f_1\) & 2 & 7 & 8 & 6 & \(\dots\)\\
\(f_2\) & 4 & 4 & 6 & 19 & \(\dots\)\\
\(f_3\) & 6 & 257 & 1 & 2 & \(\dots\)\\
\(f_4\) & 7 & 121 & 23188 & 2313 & \(\dots\)\\
\(\vdots\) & \(\vdots\) & \(\vdots\) & \(\vdots\) & \(\vdots\) & \(\ddots\)\\
\end{tabular}
\end{center}
We can now define a function \(g(n)=f_n(n)+1\) this function is total and not
in the list because it is different from any function in the list for at
least one input. As an example of such a function, we could try to define
division with rest on natural numbers like the following:

\begin{minted}[fontsize=\scriptsize]{agda}
_/_ :  ℕ → ℕ → ℕ
zero / y = zero
suc x / y = suc ( (x - y) / y)
\end{minted}

The problem with this definition is that Agda doesn't know that $\mathtt{x-y}$\;is
smaller than $\mathtt{x+1}$, which is clearly the case ($\mathtt{x}$\;and $\mathtt{y}$\;are positive).  This
definition would work perfectly fine in a language without termination
checking (like Haskell).  Agda only checks if an argument is structurally
decreasing.  Here, it is neither the case for $\mathtt{x}$\;nor for $\mathtt{y}$.

To remedy this problem sized types were introduced first to Mini-Agda (a
language specifically developed to explore them) by \cite{abel2010miniagda}.
Later, they got introduced to Agda itself. Sized types allow us to annotate data
with their size. Functions can use these sizes to check termination and
productivity.

We can now define the natural numbers depending on a size argument.
\begin{minted}[fontsize=\scriptsize]{agda}
data ℕ (i : Size) : Set where
  zero : ℕ i
  suc : ∀{j : Size< i} → ℕ j → ℕ i
\end{minted}
The natural number now depends on the size $\mathtt{i}$.  The constructor $\mathtt{zero}$\;is of
arbitrary size $\mathtt{i}$.  The constructor $\mathtt{suc}$\;gets a size $\mathtt{j}$\;which is smaller than $\mathtt{i}$, a
natural number of size $\mathtt{j}$\;and gives back a natural number of size $\mathtt{i}$.  This
means the size of the input is smaller than the size of the output.  For
inductive types, size is an upper bound on the number of constructors.  With
$\mathtt{suc}$\;we add a constructor so the size has to increase  $\mathtt{i}$.  We can now define
subtraction on these sized natural numbers.
\begin{minted}[fontsize=\scriptsize]{agda}
_-_ : {i : Size} → ℕ i → ℕ ∞ → ℕ i
zero    - _      = zero
n       - zero   = n
(suc n) - (suc m) = n - m
\end{minted}
Through the sized annotations we know now that the result isn't larger than
the first input.  \(\infty\) means that the size isn't bound.  If the first
argument is zero the result is also zero, which has the same type.  If the
second argument is zero we return just the first.  In the last, case both
arguments are non-zero.  We call subtraction recursively on the predecessors
of the inputs.  Here, the size and both arguments are smaller.  So the
function terminates.  Though the type is smaller than \(i\), the result type
checks because sizes are upper bounds.  We can now define division.
\begin{minted}[fontsize=\scriptsize]{agda}
_/_ : {i : Size} → ℕ i → ℕ ∞ → ℕ i
zero  / _ = zero
suc x / y = suc ( (x - y) / y)
\end{minted}
From the definition of $\mathtt{suc}$\;we know that the size of $\mathtt{x}$\;is smaller than $\mathtt{i}$.
Because the result of $\mathtt{-}$\;has the same size as its first input (here $\mathtt{x}$),
we also know that $\mathtt{(x\;-\;y)}$\;has the same size as $\mathtt{x}$. Therefore, $\mathtt{(x\;-\;y)}$\;is
smaller than $\mathtt{suc\;x}$\;and the function is decreasing on the first argument.
Also, Agda accepts this definition.

We can also use sized types for coinductive types.  To show this we will
define the hamming function.  This produces a stream of all composites of two
and three in order.  First, we will define the sized stream type.
\begin{minted}[fontsize=\scriptsize]{agda}
record Stream (i : Size) (A : Set) : Set where
  coinductive
  field
    hd : A
    tl : ∀ {j : Size< i} → Stream j A
open Stream
\end{minted}
This stream has a new parameter of type $\mathtt{Size}$. This size gives the minimal
definition depth of the stream. The definition depth says how often we can
destruct the stream without diverging. If we take the tail of a stream, the
output stream's depth would be one smaller. Because in Agda coinductive types
can't have indexes, we can only say that its depth is smaller. We will now
define some helper functions for the hamming function. First, we need a cons
function.
\begin{minted}[fontsize=\scriptsize]{agda}
cons : {i : Size} {A : Set} → A -> Stream i A → Stream i A
hd (cons x _)  = x
tl (cons _ xs) = xs
\end{minted}
This just appends an element at the front of the stream.  Because the output
stream's depth is larger than the input and the size is a minimum, we can give
the output the same size parameter as the input.  Now we will define map over
streams.
\begin{minted}[fontsize=\scriptsize]{agda}
map : {A B : Set} {i : Size} → (A → B) → Stream i A → Stream i B
hd (map f xs) = f (hd xs)
tl (map f xs) = map f (tl xs)
\end{minted}
This function just changes the content of the stream so the size stays the
same.  The last helper function we need is the merge function.
\begin{minted}[fontsize=\scriptsize]{agda}
merge : {i : Size} → Stream i ℕ → Stream i ℕ → Stream i ℕ
hd (merge xs ys) = hd xs ⊓ hd ys
tl (merge xs ys) = if ⌊ hd xs ≤? hd ys ⌋
                   then cons (hd ys) (merge (tl xs) (tl ys))
                   else cons (hd xs) (merge (tl xs) (tl ys))
\end{minted}
This function just merges two streams. It always compares one element of
each stream with each other and puts the bigger after the smaller. This is
clear in the case for $\mathtt{hd}$\;($\mathtt{\sqcup}$\;is just the binary minimum function in
Agda). In the $\mathtt{tl}$\;case we just compare the heads of the stream and construct
the tail with $\mathtt{cons}$\;accordingly. Both input streams have a minimal
definition depth of $\mathtt{i}$. Because $\mathtt{cons}$\;isn't destructing the stream (the
minimal depth doesn't get smaller) we can say that the minimum depth of the
output also won't get smaller. With all this function we can now define the
ham function. Here it is:
\begin{minted}[fontsize=\scriptsize]{agda}
ham : {i : Size} → Stream i ℕ
hd ham = 1
tl ham = (merge (map (_*_ 2) ham) (map (_*_ 3) ham))
\end{minted}
None of the used function is destructing the stream, so this definition gets
accepted.
\chapter{Type Theory based on dependent Inductive and Coinductive Types}
\label{sec:org33f714d}
In the paper \cite{basold2016type} a type theory, where inductive types and
coinductive types can depend on values, is developed. For example, we can, in
contrast to the coinductive types of Coq and Agda, define streams which depend
on their definition length. The theory differentiates types from terms. We
don't have infinite universes, where a term in universe \(n\) has a type in
universe \(n+1\)(This is how it is done in Coq \cite{sozeau2014universe} and Agda
\cite{agdadocuniverselevels}). Therefore, types can only depend on values, not on
other types. We only have functions on the type level. These functions abstract
over terms. For example, \(\lambda x.A\) is a type where all occurrences of the
term variable \(x\) in \(A\) are bound.  We will see that functions are definable
on the term level. We can apply types to terms. For example, \(A @ t\) means we
apply the term \(A\) to \(x\). Every type has a kind. A kind is either \(*\) or
\(\Gamma\rat*\). Here, \(\Gamma\) is a context which states to what terms we can
apply the type. For example, we can apply \(A\) of kind \((x:B)\rat*\) only to a
term of type \(B\). If we apply it to \(t\) of type \(B\), we get a type of kind
\(*\). We write \(\rat\) instead of \(\rightarrow\) to indicate, that these are not
functions. We can also apply a term to another term. For example, \(t@s\) means
we apply the term \(t\) to the term \(s\).  Terms also can depend on contexts.
For example, if we have a term \(t\) of type \((x:A)\rat B\) and apply it to a term
\(s\) of type \(A\) we get a term of type \(B\).  We can also define our own types.
\(\mu(X:\Gamma\rat*;\vv{\sigma};\vv{A})\) is an inductive type and
\(\nu(X:\Gamma\rat*;\vv{\sigma};\vv{A})\) is a coinductive type. \(X\) is a
variable that stands for the recursive occurrence of the type. It has the
same kind \(\Gamma\rat*\) as the defined type. The \(\vv{A}\) can contain this
variable. There are also contexts \(\vv{\Gamma}\), which are implicit in the
paper. \(\sigma_k\) and \(A_k\) can contain variables from \(\Gamma_k\). \(\sigma_k\)
is a context morphism from \(\Gamma_k\) to \(\Gamma\). A context morphism is a
sequence of terms, which depend on \(\Gamma_k\) and instantiate \(\Gamma\).
\(\vv{\sigma}\), \(\vv{A}\) and \(\vv{\Gamma}\) are of the same length.

In this theory, we can define partial streams on some type \(A\) like the following:
\begin{align*}
&\text{PStr }A := \nu(X:(n:\text{Conat})\rat*;(\text{succ} @ n, \text{succ} @ n);(A, X @ n))\\
&\text{with } \Gamma_1 = (n:\text{Conat}) \text{ and } \Gamma_2 = (n:\text{Conat})
\end{align*}
Here, $\mathtt{succ}$\;is the successor on co-natural numbers.  Co-natural numbers are
natural numbers with one additional element, infinity. See \ref{sec:orgc62542e}
for their definition. Here, the first destructor is the head. It becomes a
stream with length \(\text{succ} @ N\) and returns an \(A\). The second destructor
is the tail. It becomes also a stream of length \(\text{succ} @ N\). It gives
back an \(X @ n\), which is a stream of length \(n\). We can also define the Pi
type from \(A\) to \(B\), where \(B\) can depend on \(A\).
\begin{align*}
&\Pi x:A.B := \nu(\_:*;\epsilon_1;B)\\
&\text{with } \Gamma_1 = (x:A)
\end{align*}
By \(\_\) we mean, we are ignoring this variable. \(\epsilon_1\) is one empty
context morphism.  So the only destructor gives back a \(B\) which can depend
on \(x\) of type \(A\).  It is the function application.

To construct an inductive types we use constructors (written
\(\alpha_k^{\mu(X:\Gamma\rat*;\vv{\sigma};\vv{A})}\) in the paper, which is the k-th
constructor of the given type).  We can destruct it with recursion (written
$\mathtt{rec}$\;\(\vv{(\Gamma_k.y_k).g_k}\)).  Coinductive type work the other way around.
We destruct them with destructors (written
\(\xi_k^{\nu(X:\Gamma\rat*;\vv{\sigma};\vv{A})}\)) and construct them with
corecursion (written $\mathtt{corec}$\;\(\vv{(\Gamma_k.y_k).g_k}\)).

We will give the rules for the theory in Section \ref{sec:orga5d1bdb} and a detailed
explanation of the reduction in \ref{sec:orgf9c7a7b}.

\chapter{Implementation}
\label{sec:org8542692}
In this section, we look at the implementation details.  We use the functional
programming language Haskell for implementing the theory. Haskell is a pure
language. This means functions which aren't in the IO monad have no side
effects. The only IO we are doing is reading a file and as the last step
printing it. Because everything between this is pure, we can test it without
bordering on side effects. Another feature of Haskell, which will get
useful in our implementation is pattern matching. We will see its usefulness
in Section \ref{sec:orga5d1bdb}.

In Section \ref{sec:org9120160} we will develop the abstract syntax of our language
from the raw syntax in the paper. Then, we rewrite the typing rules in \ref{sec:orga5d1bdb}. At last we look at the implementation of the reduction in \ref{sec:orgf9c7a7b}
\section{Abstract Syntax}
\label{sec:org9120160}
In the following, we will scratch out the abstract syntax. In contrast to
\cite{basold2016type} we can't write anonymous inductive and coinductive types.
We will give every inductive and coinductive type a name. They will be
defined via declarations. In these declarations, we will give, their
constructors/destructors. They will also be given names. In \cite{basold2016type}
they are anonymous. We can then refer to the previously defined types. We
will describe declarations in Section \ref{sec:orgf762281}. We will also be able to
bind expressions to names. In Section \ref{sec:orgdbcd04a} we will define the syntax
of expressions. This will mostly be in one to one correspondence with the
syntax of \cite{basold2016type}. Note however, that we use the names of the
constructors instead of anonymous constructors together with their type and
number. Also, the order of the matches in $\mathtt{rec}$\;and $\mathtt{corec}$\;is irrelevant. We use
the names of the Con/Destructors to identify them. In the following Section
\ref{sec:orgc5acffd}, we will see how the examples from the paper look in our concrete
syntax.
\subsection{Declarations}
\label{sec:orgf762281}
The abstract syntax is given in Figure \ref{syntax-for-declarations}. With the
keywords data and codata we define inductive and coinductive types
respectively. After that, we will write the name. We can only use names that
aren't used already. Behind that, we can give a parameter context. This is a
type context. These types are not polymorphic. They are merely macros to make
the code more readable and allow the definition of nested types. If we want
to use these types we have to fully instantiate this context. These types can
occur everywhere in the definition where a type is expected. A (co)inductive
type can have a context which is written before an arrow. $\mathtt{Set}$\;stands for
type (or * in the paper). If a type doesn't have a context we omit the arrow.
We will also give names to every constructor and destructor. These names have
to be unique. Constructors and destructors also have contexts. Additionally,
they have one argument which can have a recursive occurrence of the type we
are defining. A constructor gives back a value of the type, where its
context is instantiated. This instantiation corresponds to the sigmas in the
paper. If we write a name before an equal sign we can bind the following
expression to the name. Every such defined name can depend on a parameter
context and an argument context. We write the parameter context like in the
case for data types behind the name. After that, we can give a term context
between round parenthesis.

\begin{figure}
 \begin{align*}
   \begin{array}{llll}
    N &:= &[A-Z][a-zA-Z0-9]* &\text{Names for types,}\\
         & &&          \text{constructors}\\
         & &&          \text{and destructors}\\
    n &:= &[a-z][a-zA-Z0-9]* &\text{Names for expressions}\\
    EV &:= &x,y,z,\dots &\text{Expression variables} \\
    TV &:= &X,Y,Z,\dots &\text{Type expression} \\
         & &&          \text{variables}\\
    PV  &:= &A,B,C,\dots &\text{Parameter variables} \\
    EC &:= &\lozenge &\text{Expression Context} \\
           &| &\text{(} EV \text { : } TV (,EV \text{ : } TV)*\text{)}& \\
    PC &:= &\langle\rangle &\text{Parameter Context} \\
       &| &\langle(PV \text{ : } EC \rightarrow \text{ Set})*\rangle & \\
    Decl &:= &\text{data } N\; PC \text{ : } (EC \rightarrow)? \text{ Set where} &\text{Declarations}\\
              &&\quad(N \text{ : }  (EC \rightarrow)? TypeExpr \rightarrow N\; Expr*)* &\\
         &| &\text{codata } N\; PC \text{ : } (EC \rightarrow)? \text{ Set where}& \\
     &&\quad(N \text{ : }  (EC \rightarrow)? N\; Expr* \rightarrow TypeExpr)*& \\
    &| &n \; PC \; EC \text{ = } Expr & \\
  \end{array}
\end{align*}
\caption{Syntax for declarations}
\label{syntax-for-declarations}
\end{figure}

The declarations in Figure \ref{syntax-for-declarations} correspond to \(\rho(X:\Gamma\rat*;\vv\sigma;\vv{A}):\Gamma\rat*\) as follows:
\begin{itemize}
\item The first \(N\) is X
\item The other \(N\) will be used later for
\(\alpha_1^{\mu(X:\Gamma\rat *;\vv\sigma;\vv A)},\alpha_2^{\mu(X:\Gamma\rat *;\vv\sigma;\vv A)},\dots\)
in the case of inductive types and
\(\xi_1^{\nu(X:\Gamma\rat *;\vv\sigma;\vv A)},\xi_2^{\nu(X:\Gamma\rat *;\vv\sigma;\vv A)},\dots\)
in the coinductive case
\item The \(TypExpr\) are the \(\vv{A}\)
\item The \(Expr*\) are the \(\vv{\sigma}\)
\item The first \(EC\) is \(\Gamma\)
\item The other \(EC\) stand for \(\Gamma_1,\dots,\Gamma_m\)
\end{itemize}

To parse the abstract syntax we use Megaparsec. The parser generates an
abstract syntax tree, which is given for declarations in Listing
\ref{Abstract Syntax Tree for Declarations}. The field $\mathtt{ty}$\;in $\mathtt{ExprDef}$\;is used later in
type checking. The parser just fills them in with $\mathtt{Nothing}$. Data and codata
definitions are both saved in $\mathtt{TypeDef}$. The Haskell type $\mathtt{OpenDuctive}$\;contains all the
information for inductive and coinductive types. It corresponds to \(\mu\) and
\(\nu\) in the paper. We use an $\mathtt{OpenDuctive}$\;where the field $\mathtt{inOrCoin}$\;is $\mathtt{IsIn}$
for \(\mu\) and an $\mathtt{OpenDuctive}$\;where the field $\mathtt{inOrCoin}$\;is $\mathtt{IsCoin}$\;for
\(\nu\).  The Haskell type $\mathtt{StrDef}$\;ensures that the sigmas, as and gamma1s have the
same length.  We omit the implementation details for the parser because we
are mainly focused on type checking.

\begin{listing}[htbp]
\begin{minted}[fontsize=\scriptsize]{haskell}
data Decl = ExprDef { name :: Text
                    , tyParameterCtx :: TyCtx
                    , exprParameterCtx :: Ctx
                    , expr :: Expr
                    , ty :: Maybe Type
                    }
          | TypeDef OpenDuctive
          | Expression Expr

data OpenDuctive = OpenDuctive { nameDuc :: Text
                               , inOrCoin :: InOrCoin
                               , parameterCtx :: TyCtx
                               , gamma :: Ctx
                               , strDefs :: [StrDef]
                               }

data StrDef = StrDef { sigma :: [Expr]
                     , a :: TypeExpr
                     , gamma1 :: Ctx
                     , strName :: Text
                     }
\end{minted}
\caption{\label{Abstract Syntax Tree for Declarations}
Implementation of the abstract syntax of fig. \ref{syntax-for-declarations}}
\end{listing}
\subsection{Expressions}
\label{sec:orgdbcd04a}
The abstract syntax for expression is given in Figure \ref{syntax-for-expressions}.
We will separate expressions in expressions for terms and expressions for
types.  There are given as regular expressions in $\mathtt{Expr}$\;and $\mathtt{TypeExpr}$\;respectively.

\begin{figure}
 \begin{align*}
   \begin{array}{llll}
     ParInst &:= &\langle TypeExpr(\text{,}TypeExpr)*\rangle &\text{Instantiations for}\\
                                                            &&&\text{paramter contexts}\\
     ExprInst &:= &\text{(}Expr(\text{,}Expr)*\text{)} &\text{Instantiations for}\\
                                                       &&&\text{expression contexts}\\
     Expr &:= &\text{rec } N \; ParInst? \text{ to } TypeExpr \text{ where} &\text{expression}\\
     &&\quad Match*&\\
     &| &\text{corec } TypeExpr \text { to } N \; ParInst? \text{ where}&\\
     &&\quad Match*&\\
     &| &Expr\text{ @ }Expr &\\
     &| &\lozenge&\\
     &| &EV&\\
     &| &n\; ParInst\; ExprInst&\\
     Match &:= &N\; EV* = Expr &\text{match}\\
     TypeExpr &:= &\text{(}EV\text{ : }TypeExpr\text{).}TypeExpr &\text{Type expressions}\\
     &| &TypeExpr\text{ @ }Expr&\\
     &| &\text{Unit} &\\
     &| &TV&\\
     &| &N\; ParInst? &\\
  \end{array}
\end{align*}
\begin{lstlisting}
\end{lstlisting}
\caption{Syntax for expressions}
\label{syntax-for-expressions}
\end{figure}

An $\mathtt{Expr}$\;is either a $\mathtt{rec}$, a $\mathtt{corec}$, a con/destructor, an application
$\mathtt{@}$, the only primitive unit expression $\mathtt{\lozenge}$\;or a variable. With the
keyword $\mathtt{rec}$\;we can destruct an inductive type. We write
$\mathtt{N\;ParInst?\;to\;TypeExrp}$, where $\mathtt{N}$\;is a previously defined inductive type
and $\mathtt{ParInst?}$\;the instantiation of its parameter context, after $\mathtt{rec}$\;to
facilitate type checking. It says we want to destruct an inductive type to
some other type. We have to list all the constructors above one another.
For each constructor, we write an expression behind the equal sign, which
should be of type $\mathtt{TypeExpr}$\;which we have given above. In this expression,
we can use variables given in the match expression. The last one is the
recursive occurrence. With the keyword $\mathtt{corec}$\;we can do the same thing to
construct a coinductive type. Here, we have to swap the $\mathtt{N\;ParInst?}$\;and the
$\mathtt{TypeExpr}$\;and list the destructors. All con/destructors have to be
instantiated with all variables in the parameter contexts of their types.
This is done by giving types of the expected kinds separated by ',' enclosed
in $\mathtt{\langle}$\;and $\mathtt{\rangle}$. The variables are separated into local variables
and global variables. Global variables refer to previously defined
expressions. We have to fully instantiate their parameter contexts and their
expression contexts. We can also apply an expression to another with $\mathtt{@}$.
This application is left-associative.  So if we write $\mathtt{t\;@\;s\;@\;v}$\;we mean
$\mathtt{(t\;@\;s)\;@\;v}$.

The $\mathtt{typeExpr}$\;is either the unit type $\mathtt{Unit}$, a lambda abstraction on
types, an application, or a variable. In the lambda expression, we have to
give the type of the variable. We apply a type to a term (types can only
depend on terms) with $\mathtt{@}$.  As in the case of term application, this is
also left-associative.  The unit type is the only primitive type
expression.

The generated abstract syntax tree is given in Listing
\ref{abstract-syntax-tree-for-expressions}. The variables for expressions are
separated in $\mathtt{LocalExprVar}$\;and $\mathtt{GlobalExprVar}$. $\mathtt{LocalExprVar}$\;should refer
to variables that are only locally defined i.e. in $\mathtt{Rec}$\;and $\mathtt{Corec}$. We
use de Bruijn indexes for them. This facilitates substitution which we will
describe in Section \ref{sec:orgd528f7c}. $\mathtt{GlobalExprVar}$\;refers to variables from
definitions. Here, we just use names. We do the same thing for $\mathtt{LocalTypeVar}$
and $\mathtt{GlobalTypeVar}$. In the abstract syntax tree, we use anonymous
constructors like in the paper. We combine them with the Haskell constructor
$\mathtt{Structor}$. We know from the field $\mathtt{ductive}$\;if it is a constructor or a
destructor. The types in field $\mathtt{parameters}$\;are to fill in the parameter
context of the field $\mathtt{ductive}$. The field $\mathtt{nameStr}$\;in $\mathtt{Constructor}$\;and
$\mathtt{Destructor}$\;are just for printing. We combine $\mathtt{rec}$\;and $\mathtt{corec}$\;to $\mathtt{Iter}$.

\begin{listing}[htbp]
\begin{minted}[fontsize=\scriptsize]{haskell}
data TypeExpr = UnitType
              | TypeExpr :@ Expr
              | LocalTypeVar Int Bool Text
              | Parameter Int Bool Text
              | GlobalTypeVar Text [TypeExpr]
              | Abstr Text TypeExpr TypeExpr
              | Ductive { openDuctive :: OpenDuctive
                        , parametersTyExpr :: [TypeExpr]}

data Expr = UnitExpr
          | LocalExprVar Int Bool Text
          | GlobalExprVar Text [TypeExpr] [Expr]
          | Expr :@: Expr
          | Structor { ductive :: OpenDuctive
                     , parameters :: [TypeExpr]
                     , num :: Int
                     }
          | Iter { ductive :: OpenDuctive
                 , parameters :: [TypeExpr]
                 , motive :: TypeExpr
                 , matches :: [([Text],Expr)]
                 }
\end{minted}
\caption{\label{abstract-syntax-tree-for-expressions}
Implementation of the abstract syntax of fig. \ref{syntax-for-expressions}}
\end{listing}

\section{Substitution}
\label{sec:orgd528f7c}
In the following we will write \(t[s/x]\) for "substitute every free
occurrences of \(x\) in \(t\) by \(s\)". Substitution is done in the module
$\mathtt{Subst.hs}$. We use de Bruijn indexes \cite{de1972lambda} for bound variables to facilitate
substitution. With this method, every bound variable is a number instead of a
string. The number says where the variable is bound. To find the binder of a
variable we go outwards from it and count every binder until we reach the
number of the variable. For example, \(\lambda.\lambda.\lambda.1\) says that the
variable is bound by the second binder (we start counting at zero). This
would be the same as \(\lambda x.\lambda y. \lambda z.y\). This means we never
have to generate fresh names. We just shift the free variables in the term
with which we substitute by one, every time we encounter a binder. This
shifting is done in the module $\mathtt{ShiftFreeVars.hs}$. We also want to be able to
substitute multiple variables simultaneously. If we would just substitute one
term after another we could substitute into a previous term. For example, the
substitution \(x[y/x][z/y]\) would yield \(z\) if we substitute sequential and
\(y\) if we substitute simultaneously.  To make simultaneous substitution
possible every local variable has a boolean flag.  If this flag is set to
true substitution won't substitute for that variable.  So for simultaneous
substitutions, we just set this flag to true for all terms with which we want
to substitute.  Then, we substitute with them.  In the last step, we just have
to set the flags to false in the result.  This setting(marking of the
variables) is done in the module $\mathtt{Mark.hs}$.

\section{Typing rules}
\label{sec:orga5d1bdb}
A typing rule says that some expression or declaration is of some type, given
some premises. If we can for every declaration or expression form a tree of
such rules with no open premises, our program type checks. We have to rewrite
the typing rules of the paper, to get rules which are syntax-directed.
Syntax-directed means we can infer from the syntax alone what we have to check next
i.e. which rule with which premises we have to apply. In the paper, there are
rules containing variables in the premises where their type isn't in the
conclusion. So if we want to type-check something which is the conclusion of
such a rule we have no way of knowing what these variables are.

We don't need the weakening rules because we can look up a variable in a
context.  So we ignore them in our implementation.

The order in \textbf{TyCtx} isn't relevant so we can use a map for it. In the code, we
use a list because the names of the variables are the index of their type in
the context. The order of \textbf{Ctx} is relevant because types of later variables
can refer to former variables and application instantiates the first variable
in \textbf{Ctx}. We add a new context for data types. We also need a context for the
parameters. \textbf{Ctx} can contain variables from this context, but not from
\textbf{TyCtx}.

We also rewrite the rules which are already syntax-directed to rules which
work on our syntax.   We will mark semantic differences in the rewritten rules
gray. We use variables \(\Phi,\Phi',\Phi_1,\Phi_2,\dots\) for parameter contexts,
\(\Theta,\Theta',\Theta_1,\Theta_2,\dots\) for type variable contexts and
\(\Gamma,\Gamma',\Gamma_1,\Gamma_2,\dots\) for term variable contexts.
The judgments in our rules are of one of the following form.
\begin{itemize}
\item \(\Phi\mid\Theta\mid\Gamma\vdash\Theta'\) - The type variable context
\(\Theta'\) is well-formed in the combined context \(\Phi\mid\Theta\mid\Gamma\).
\item \(\Phi\mid\Theta\mid\Gamma\vdash\Gamma'\) - The term variable context
\(\Gamma'\) is well-formed in the combined context \(\Phi\mid\Theta\mid\Gamma\).
\item \(\Phi\mid\Theta\mid\Gamma\vdash\Phi'\) - The parameter variable context
\(\Phi'\) is well-formed in the combined context \(\Phi\mid\Theta\mid\Gamma\).
\item \(A\longrightarrow_T^* B\) - The type \(A\) fully evaluates to type \(B\).
\item \(A \equiv_\beta B\) - The type \(A\) is computational equivalent to type \(B\).
\item \(\Phi\mid\Theta\mid\Gamma\vdash A : \Gamma_2\rat*\) - The type
\(A\) is well-formed in the combined context \(\Phi\mid\Theta\mid\Gamma\) and
can be instantiated with arguments according to context \(\Gamma_2\).
\item \(\Phi\mid\Theta\mid\Gamma\vdash t : \Gamma_2\rat A\) - The term \(t\) is
well-formed in the combined context \(\Phi\mid\Theta\mid\Gamma\) and can be
instantiated with arguments according to context \(\Gamma_2\).  After this
instantiation, it is of type \(A\), where the arguments are substituted in \(A\).
\item \(\Phi \vdash \sigma : \Gamma_1 \triangleright \Gamma_2\) - The context
morphism \(\sigma\) is a well-formed substitution for \(\Gamma_2\) with terms
in context \(\Gamma_1\) in parameter context \(\Phi\).
\end{itemize}
We will write \(\vdash\) for \(\Phi\mid\Theta\mid\Gamma\vdash\) where
\(\Phi\),\(\Theta\), and \(\Gamma\) are arbitrary and aren't referred to by the
right-hand side.

In the module $\mathtt{TypeChecker}$\;we will implement the following rules.  It
defines a monad $\mathtt{TI}$\;which can throw errors and has a reader on the contexts
in which we are type checking.  To add something to a context we use the
function $\mathtt{local}$.  This function gets a function to change the current
content of the reader monad and executes a reader on this changed context in
the current monad.

\subsection{Context rules}
\label{sec:orgb4333f6}
The rules for valid contexts are already-syntax directed so we take
just them.
\begin{center}
\AxiomC{}
\UnaryInfC{$\vdash\emptyset$ \TyCtx}
\DisplayProof
\hskip 1.5em
\AxiomC{$\vdash\Theta$ \TyCtx}
\AxiomC{$\vdash\Gamma$ \Ctx}
\BinaryInfC{$\vdash\Theta,X:\Gamma\rat*$ \TyCtx}
\DisplayProof
\vskip 0.5em
\AxiomC{}
\UnaryInfC{$\vdash\emptyset$ \Ctx}
\DisplayProof
\hskip 1.5em
\AxiomC{$\mid\emptyset\mid\Gamma\vdash A:*$}
\UnaryInfC{$\vdash\Gamma,x:A$ \Ctx}
\DisplayProof
\end{center}
In the rules for valid contexts, we ensure that the types in the context can
not depend on \textbf{TyCtx}.  Note however that they can depend on \textbf{ParCtx}.  This
ensures that only strictly positive types are possible.

We also need new rules for checking if a parameter context is valid.
\begin{center}
\AxiomC{}
\UnaryInfC{$\vdash\emptyset$ \ParCtx}
\DisplayProof
\hskip 1.5em
\AxiomC{$\vdash\Phi$ \ParCtx}
\AxiomC{$\vdash\Gamma$ \Ctx}
\BinaryInfC{$\vdash\Phi,X:\Gamma\rat*$ \ParCtx}
\DisplayProof
\end{center}
These are structural the same rules like this for \textbf{TyCtx}.  The difference is that \textbf{ParCtx}
and \textbf{TyCtx} are used differently in the other rules, as we have already seen
in the rule for \textbf{Ctx}.

We use the notation \(\Theta(X)\rightsquigarrow\Gamma\rat*\) for looking up
the type variable \(X\) in type context \(\Theta\) yields type \(\Gamma\rat*\). We
add 2 rules for looking up something in a type context. They are:
\begin{center}
  \AxiomC{$\vdash \Theta$ \TyCtx}
  \AxiomC{$\vdash \Gamma$ \Ctx}
  \BinaryInfC{$\Theta,X:\Gamma\rat*(X)\rightsquigarrow\Gamma\rat*$}
  \DisplayProof
  \hskip 1.5em
  \AxiomC{$\vdash \Gamma_1$ \Ctx}
  \AxiomC{$\Theta(X) \rightsquigarrow\Gamma_2\rat*$}
  \BinaryInfC{$\Theta,Y:\Gamma_1\rat*(X)\rightsquigarrow\Gamma_2\rat*$}
  \DisplayProof
\end{center}
Here, \(Y\) and \(X\) are different variables.

The rules for looking up something in a parameter context are principally the
same.
\begin{center}
  \AxiomC{$\vdash \Phi$ \ParCtx}
  \AxiomC{$\vdash \Gamma$ \Ctx}
  \BinaryInfC{$\Phi,X:\Gamma\rat*(X)\rightsquigarrow\Gamma\rat*$}
  \DisplayProof
  \hskip 1.5em
  \AxiomC{$\vdash \Gamma_1$ \Ctx}
  \AxiomC{$\Phi(X) \rightsquigarrow\Gamma_2\rat*$}
  \BinaryInfC{$\Phi,Y:\Gamma_1\rat*(X)\rightsquigarrow\Gamma_2\rat*$}
  \DisplayProof
\end{center}

Respectively the notation \(\Gamma(x)\rightsquigarrow A\) means looking
up the term variable \(x\) in term context \(\Gamma\) yields type \(A\). The
rules for term contexts are:
\begin{center}
  \AxiomC{$\vdash \Gamma$ \Ctx}
  \AxiomC{$\Gamma\vdash A:*$}
  \BinaryInfC{$\Gamma,x:A(x)\rightsquigarrow A$}
  \DisplayProof
  \hskip 1.5em
  \AxiomC{$\Gamma(x) \rightsquigarrow A$}
  \AxiomC{$\Gamma\vdash B:*$}
  \BinaryInfC{$\Gamma,y:B(x)\rightsquigarrow A$}
  \DisplayProof
\end{center}

\subsection{Beta-equivalence}
\label{sec:org6ac2816}
Two types are beta equivalent if they evaluate to the same type. Because our
language is deterministic this just means if we fully evaluate both of them
they are alpha equivalent. Alpha equivalence means we can substitute some
variables in both of them and get the same type. So we first need to define
rules which say what full evaluation means. We write \(A \longrightarrow_T^*
    B\) for evaluating \(A\) as long as it is possible yields \(B\).

The rules are:
\begin{center}
\AxiomC{$\neg\exists B : A \longrightarrow_T B$}
\UnaryInfC{$A \longrightarrow_T^* A$}
\DisplayProof
\hskip 1.5em
\AxiomC{$A \longrightarrow_T B$}
\AxiomC{$B \longrightarrow_T^* C$}
\BinaryInfC{$A \longrightarrow_T^* C$}
\DisplayProof
\end{center}
\(\longrightarrow_T\) is defined in Section \ref{sec:orgf9c7a7b}.

We can then introduce a new rule for beta-equivalence.
\begin{center}
\AxiomC{$A\longrightarrow_T^* A'$}
\AxiomC{$B\longrightarrow_T^* B'$}
\AxiomC{$A'\equiv_\alpha B'$}
\TrinaryInfC{$A\equiv_\beta B$}
\DisplayProof
\end{center}
This rule says if \(A\) evaluates to \(A'\), \(B\) to \(B'\) and \(A'\) and \(B'\) are
alpha equivalent, then \(A\) and \(B\) are beta equivalent. In the
implementation \(\equiv_\alpha\) is trivial because we use \emph{de Bruijn
indices}.

We also add some rules to check if two contexts are the same.
   \begin{center}
   \AxiomC{}
   \UnaryInfC{$\emptyset\equiv_\beta\emptyset$}
   \DisplayProof
   \hskip 1.5em
   \AxiomC{$\Gamma_1\equiv_\beta \Gamma_2$}
   \AxiomC{$A\equiv_\beta B$}
   \BinaryInfC{$\Gamma_1,x:A\equiv_\beta\Gamma_2,y:B$}
   \DisplayProof
%   \vskip 0.5em
%   \AxiomC{$\Theta_1\equiv_\beta \Theta_2$}
%   \AxiomC{$\Gamma_1\equiv_\beta \Gamma_2$}
%   \BinaryInfC{$\Theta_1,X:\Gamma_1\rat*\equiv_\beta\Theta_2,X:\Gamma_2\rat*$}
%   \DisplayProof
   \end{center}

\subsection{Unit type and expression introduction}
\label{sec:org7f5ea3e}
The paper defines one rule for the unit type and one for the unit value.
These are.
\begin{center}
  \AxiomC{}
  \RightLabel{\textbf{($\top$-I)}}
  \UnaryInfC{$\vdash\top:*$}
  \DisplayProof
  \hskip 1.5em
  \AxiomC{}
  \topI{$\vdash\lozenge:\top$}
  \DisplayProof
\end{center}
The first rule says that the type \(\top\) has always an empty context.  The
second rule says its value \(\lozenge\) is always of type \(\top\). These rules
get rewritten to.
\begin{center}
  \AxiomC{}
  \RightLabel{\textbf{(Unit-I)}}
  \UnaryInfC{\graybox{$\Phi\mid\Theta\mid\Gamma$}$\vdash$Unit:$*$}
  \DisplayProof
  \hskip 1.5em
  \AxiomC{}
  \topI{\graybox{$\Phi\mid\Theta\mid\Gamma$}$\vdash\lozenge$:Unit}
  \DisplayProof
\end{center}
We change the syntax "\(\top\)" to "Unit" and add the contexts \(\Phi\),
\(\Theta\), \(\Gamma\). We will do this for every rule which has empty contexts
to subsume the weakening rules of the paper. The unit term always has the
unit type as its type.

\subsection{Variable lookup}
\label{sec:orgcfd2702}
We have three kinds of variables we can lookup. They are type variables,
term variables, and parameters.  The paper already has rules for the type and
term variables.  We need to rewrite them.  We add a new rule for looking up
a parameter.

The rule:
 \begin{prooftree}
  \AxiomC{$\vdash \Theta$ \TyCtx}
  \AxiomC{$\vdash \Gamma$ \Ctx}
  \TyVarI{$\Theta,X:\Gamma\rat*\mid\emptyset\vdash X : \Gamma \rat *$}
\end{prooftree}
gets rewritten to:
\begin{prooftree}
  \AxiomC{\graybox{$\Theta(X)\rightsquigarrow\Gamma\rat*$}}
  \AxiomC{\graybox{$\vdash \Gamma_1$ \Ctx}}
  \TyVarI{\graybox{$\Phi$}$\mid\Theta\mid$\graybox{$\Gamma_1$}$\vdash X : \Gamma \rat *$}
\end{prooftree}
The rule:
\begin{center}
  \AxiomC{$\Gamma\vdash A:*$}
  \RightLabel{\textbf{(Proj)}}
  \UnaryInfC{$\Gamma,x:A\vdash x:A$}
  \DisplayProof
\end{center}
gets rewritten to:
\begin{center}
  \AxiomC{\graybox{$\Gamma(x)\rightsquigarrow A$}}
  \RightLabel{\textbf{(Proj)}}
  \UnaryInfC{\graybox{$\Phi\mid\Theta\mid$}$\Gamma\vdash x:A$}
  \DisplayProof
\end{center}
The rule for looking something up in the parameter context is:
\begin{prooftree}
  \AxiomC{$\Phi(X)\rightsquigarrow\Gamma\rat*$}
  \AxiomC{$\vdash \Gamma_1$ \Ctx}
  \TyVarI{$\Phi\mid\Theta\mid$$\Gamma_1$$\vdash X : \Gamma \rat *$}
\end{prooftree}

In the rule from the paper, we can only infer the type or kind of the last variable in the
context.  In our rules, we just look up the variable in the context.  These
rules can check the same thing if we take the weakening rules into account.
With them, we can just weaken the context until we get to the desired
variable.

\subsection{Type and expression instantiation}
\label{sec:org7ba5325}
We can instantiate types and terms.  The rule:
\begin{prooftree}
  \AxiomC{$\Theta\mid\Gamma_1\vdash A:(x:B,\Gamma_2)\rat*$}
  \AxiomC{$\Gamma_1\vdash t:B$}
  \TyInst{$\Theta\mid\Gamma_1\vdash A@t:\Gamma_2[t/x]\rat*$}
\end{prooftree}
for instantiating types gets rewritten to:
 \begin{prooftree}
  \AxiomC{\graybox{$\Phi$}$\mid\Theta\mid\Gamma_1\vdash A:(x:B,\Gamma_2)\rat*$}
  \AxiomC{\graybox{$\Phi\mid\Theta$}$\mid\Gamma_1\vdash t:$\graybox{$B'$}}
  \AxiomC{\graybox{$B\equiv_\beta B'$}}
  \TyInstTrinary{\graybox{$\Phi$}$\mid\Theta\mid\Gamma_1\vdash A@t:\Gamma_2[t/x]\rat*$}
\end{prooftree}
For this rule, we have to check if \(t\) has the expected type for the first
variable in the context of \(A\).  In our version, we just infer the type for \(A\) and \(t\).
Then, we check if the first variable in the context is beta-equal to the type
of \(t\).  If that isn't the case type checking fails.  Otherwise, we just
substitute in the remaining context.

We also have a rule to instantiate terms.  This rule:
\begin{center}
  \AxiomC{$\Gamma_1\vdash t:(x:A,\Gamma_2)\rat B$}
  \AxiomC{$\Gamma_1\vdash s:A$}
  \RightLabel{\textbf{(Inst)}}
  \BinaryInfC{$\Gamma_1\vdash t@s:\Gamma_2[s/x]\rat B[s/x]$}
  \DisplayProof
\end{center}
gets rewritten to:
\begin{center}
  \AxiomC{\graybox{$\Phi\mid\Theta$}$\mid\Gamma_1\vdash t:(x:A,\Gamma_2)\rat B$}
  \AxiomC{\graybox{$\Phi\mid\Theta$}$\mid\Gamma_1\vdash s:$\graybox{$A'$}}
  \AxiomC{\graybox{$A\equiv_\beta A'$}}
  \RightLabel{\textbf{(Inst)}}
  \TrinaryInfC{\graybox{$\Phi\mid\Theta$}$\mid\Gamma_1\vdash t@s:\Gamma_2[s/x]\rat B[s/x]$}
  \DisplayProof
\end{center}
These rules are similar to the rule for type instantiation.  Here, we have to
check (or infer) a term instead of a type.  We also have to substitute \(s\) in
the result type of \(t\) (in the case of types it's always \(*\), which obviously
has no free variables).

\subsection{Parameter abstraction}
\label{sec:org451b259}
The rule:
\begin{center}
  \AxiomC{$\Theta\mid\Gamma_1,x:A\vdash B:\Gamma_2\rat*$}
  \ParamAbstr{$\Theta\mid\Gamma_1\vdash(x).B:(x:A,\Gamma_2)\rat*$}
  \DisplayProof
\end{center}
gets rewritten to:
\begin{center}
  \AxiomC{\graybox{$\Phi$}$\mid\Theta\mid\Gamma_1,x:A\vdash B:\Gamma_2\rat*$}
  \ParamAbstr{\graybox{$\Phi$}$\mid\Theta\mid\Gamma_1\vdash(x$\graybox{$:A$}$).B:(x:A,\Gamma_2)\rat*$}
  \DisplayProof
\end{center}
Here, we just add the argument of the lambda to the expression context.  Then
we check the body of the lambda.  In the syntax-directed version we have to
annotate the variable with its type, so we know which type we have to add to
the context.

\subsection{(co)inductive types}
\label{sec:orgd5ff63a}
We have to separate the rule:
\begin{prooftree}
\AxiomC{$\sigma_k:\Gamma_k\triangleright\Gamma$}
\AxiomC{$\Theta,X:\Gamma\rat*\mid\Gamma_k\vdash A_k:*$}
\FPTy
\BinaryInfC{$\Theta \mid \emptyset \vdash \rho(X : \Gamma \rat *;\vv{\sigma};\vv{A}):\Gamma\rat *$}
\end{prooftree}
into multiple rules.  First, we need rules to check the definitions of
(co)inductive types.  These are:
\begin{prooftree}
\AxiomC{$\sigma_k:\Gamma_k\triangleright\Gamma$}
\AxiomC{\graybox{$\Phi$}$\mid X:\Gamma\rat*\mid\Gamma_k\vdash A_k:*$}
\AxiomC{\graybox{$\vdash \phi$ \ParCtx}}
\FPTy
\TrinaryInfC{$\vdash$ data X$\langle\Phi\rangle$ $\Gamma \rat $ Set where; $\vv{Constr_k : \Gamma_k\rat A_k\rat X \sigma_k}$}
\end{prooftree}
\begin{prooftree}
\AxiomC{$\sigma_k:\Gamma_k\triangleright\Gamma$}
\AxiomC{\graybox{$\Phi$}$\mid X:\Gamma\rat*\mid\Gamma_k\vdash A_k:*$}
\AxiomC{\graybox{$\vdash \phi$ \ParCtx}}
\FPTy
\TrinaryInfC{$\vdash$ codata X$\langle\Phi\rangle$ : $\Gamma \rat$ Set where; $\vv{Destr_k : \Gamma_k \rat  X \sigma_k \rat  A_k}$}
\end{prooftree}
Because we only allow top-level definitions of (co)inductive types our rules
have empty contexts.  We first have to check if \(\sigma_k\) is  a context
morphism from \(\Gamma_k\) to \(\Gamma\).  This basically means that the terms
in \(\sigma_k\) are of the types in \(\Gamma\), if we check them in \(\Gamma_k\).
After that, we have to check if the \(\vv{A}\) (the arguments where we can have
a recursive occurrence) are of kind \(*\).  Because this is a top-level
definition the context \(\phi\) is provided by the code.  So we have to check
if it is valid.  We will now have to rewrite the rules for context morphism.
Here, we just add the parameter context to the rules of the paper.
\begin{center}
\AxiomC{}
\UnaryInfC{\graybox{$\Phi\vdash$}$() : \Gamma_1 \triangleright \emptyset$}
\DisplayProof
\hskip 1.5em
\AxiomC{\graybox{$\Phi\vdash$}$\sigma : \Gamma_1 \triangleright \Gamma_2$}
\AxiomC{\graybox{$\Phi\mid$}$\Gamma_1\vdash t : A[\sigma]$}
\BinaryInfC{\graybox{$\Phi\vdash$}$(\sigma,t):\Gamma_1\triangleright(\Gamma_2,x:A)$}
\DisplayProof
\end{center}
We also need a rule for the cases in which we are using these defined
variables.  This is:
\begin{prooftree}
\AxiomC{$\Phi\mid\Theta\mid\Gamma'\vdash \vv{A}:\Gamma_i \rat *$}
\UnaryInfC{$\Phi\mid\Theta\mid\Gamma'\vdash X\langle\vv{A}\rangle : \Gamma[\vv{A}]\rat *$}
\end{prooftree}
Here, \(X\) is a data or codata definition.  The parser can decide if a variable
is such a definition or a local definition. Because we are type checking
on the abstract syntax tree we also know \(\Gamma\) and \(\Phi'\). \(\Gamma\) is
just the context from the definition and \(\Phi\) is the parameter context.
Because we already typed checked this definition we just have to check if
the types given for the parameters have the right kind.  Then, we substitute
these parameters in its type.  We will now give the rules for checking if a
list of parameters matches a parameter context.
\begin{center}
\AxiomC{}
\UnaryInfC{$\Phi\mid\Theta\mid\Gamma\vdash () : ()$}
\DisplayProof
\hskip 1.5em
\AxiomC{$\Phi\mid\Theta\mid\Gamma\vdash A : \Gamma'\rat*$}
\AxiomC{$\Phi\mid\Theta\mid\Gamma\vdash \vv{A} : \Phi'[A/X]$}
\BinaryInfC{$\Phi\mid\Theta\mid\Gamma\vdash A,\vv{A} : (X:\Gamma'\rat*,\Phi'$)}
\DisplayProof
\end{center}
We just check every variable for the kinds in \(\Phi'\) one after the other.
We also have to substitute the type into the context.  Because kinds in
a parameter context can depend on variables previously defined in this context.

\subsection{Constructor and Destructor}
\label{sec:org3e1e65e}
The rule for constructors:
\begin{center}
  \AxiomC{$\mu(X:\Gamma\rat*;\vv{\sigma};\vv{A}):\Gamma\rat*$}
  \AxiomC{$1\leq k\leq\mid\vv{A}\mid$}
  \IndIBinary{$\vdash\alpha_k^{\mu(X:\Gamma\rat*;\vv{\sigma};\vv{A})}:(\Gamma_k,y:A_k[\mu/X])\rat\mu@\sigma_k$}
  \DisplayProof
\end{center}
gets rewritten to:
\begin{center}
  \AxiomC{\graybox{$\Phi\mid\Theta\mid\Gamma\vdash \vv{B} : \Phi'$}}
  \IndI{\graybox{$\Phi\mid\Theta\mid\Gamma$}$\vdash$Constr\graybox{$\langle\vv{B}\rangle$}$:(\Gamma_k\graybox{$\graybox{$[\vv{B}]$}$},y:A_k[\mu/X]\graybox{$\graybox{$[\vv{B}]$}$})\rat\mu@\sigma_k\graybox{$\graybox{$[\vv{B}]$}$}$}
  \DisplayProof
\end{center}
The rule for destructors:
\begin{center}
  \AxiomC{$\nu(X:\Gamma\rat*;\vv{\sigma};\vv{A}):\Gamma\rat*$}
  \AxiomC{$1\leq k\leq\mid\vv{A}\mid$}
  \RightLabel{\textbf{(Coind-E)}}
  \BinaryInfC{$\vdash\xi_k^{\nu(X;\Gamma\rat*;\vv{\sigma};\vv{A})}:(\Gamma_k,y:\nu@\sigma_k)\rat
    A_k[\nu/X]$}
  \DisplayProof
\end{center}
gets rewritten to:
\begin{center}
  \AxiomC{\graybox{$\Phi\mid\Theta\mid\Gamma\vdash \vv{B} : \Phi'$}}
  \RightLabel{\textbf{(Ind-I)}}
  \UnaryInfC{\graybox{$\Phi\mid\Theta\mid\Gamma$}$\vdash$Destr\graybox{$\langle\vv{B}\rangle$}$:(\Gamma_k$\graybox{$[\vv{B}]$}$,y:\nu@\sigma_k)$\graybox{$[\vv{B}]$}$\rat
    A_k[\nu/X]$\graybox{$[\vv{B}]$}$$}
  \DisplayProof
\end{center}
In the paper de/constructors are anonymous.  They come together with their
type. Therefore, we have to check if this type is valid. Constructors
construct their type. So their output value is their type \(\mu\) applied to
the context morphism \(\sigma_k\), where \(k\) is the number of the constructor.
They become as input the context \(\Gamma_k\), which is implicit in the paper,
and a value of type \(A_k[\mu/X]\), which is the type, which can contain the
recursive occurrence. Destructors are destructing their type so we get their
type \(\nu\) applied to \(\sigma_k\) as input and \(A_k[\nu/X]\) as output.

In our rules, in contrast to the paper, the de/constructors refer to some
type which we have already type-checked. We just have to check the
parameters. Every term we need is in the Haskell representation of the
de/constructor. The de/constructor has the type which we have defined in the
data definition. We just substitute the type itself for the free variable.
At last, we need to substitute the parameters for the respective variables.

\subsection{Recursion and Corecursion}
\label{sec:org316efd2}
The rule:
\begin{center}
  \AxiomC{$\vdash C:\Gamma\rat*$}
  \AxiomC{$\Delta,\Gamma_k,y_k:A_k[C/X]\vdash g_k:(C@\sigma_k)$}
  \AxiomC{$\forall k=1,\dots,n$}
  \RightLabel{\textbf{(Ind-E)}}
  \TrinaryInfC{$\Delta\vdash$ rec
    $\vv{(\Gamma_k,y_k).g_k}:(\Gamma,y:\mu@id_\Gamma)\rat C@id_\Gamma$}
  \DisplayProof
\end{center}
gets rewritten to:
\begin{scprooftree}{0.95}
  \AxiomC{$\vdash C:\Gamma\rat*$}
  \AxiomC{\graybox{$\vdash\Gamma\equiv_\beta \Gamma'[\vv{D}]$}}
  \noLine
  \UnaryInfC{\graybox{$\vv{\vdash B_k\equiv_\beta(C@\sigma_k[\vv{D}])}$}}
  \AxiomC{\graybox{$\Phi\mid\Theta\mid\Delta\vdash \vv{D}:\Phi'$}}
  \noLine
  \UnaryInfC{$\vv{$\graybox{$\Phi\mid\mid$}$\Delta,\Gamma_k$\graybox{$[\vv{D}]$}$,y_k:A_k$\graybox{$[\vv{D}]$}$[C/X]\vdash g_k:\text{\graybox{$B_k$}}}$}
  \RightLabel{\textbf{(Ind-E)}}
  \TrinaryInfC{\graybox{$\Phi\mid\Theta\mid$}$\Delta\vdash$ rec \graybox{$\mu\langle\vv{D}\rangle$ to C};
    $\vv{\text{Constr}_k\vv{x_k}\text{ } y_k = g_k}:(\Gamma,y:\mu$\graybox{$[\vv{D}]$}$@id_\Gamma)\rat C@id_\Gamma$}
 \end{scprooftree}

We are recursing over some previously inductively defined type \(\mu\) to some
type \(C\).  These types must have the same context.  Recursing is done by
Listing each constructor with the result, which the whole expression should
have if we apply it to this constructor.  This result can refer to the
arguments of the constructor via the variables \(\vv{x_k},y_k\).  The type
must be the result type \(C\) applied to the \(\sigma_k\) of this constructor.
In the syntax-directed version, we also have to check the parameters.  We
check if the types match by inferring them and compare them on beta
equality.

We have a similar rule for corecursion.  It:
\begin{center}
  \AxiomC{$\vdash C:\Gamma\rat*$}
  \AxiomC{$\Delta,\Gamma_k,y_k:(C@\sigma_k)\vdash g_k:A_k[C/X]$}
  \AxiomC{$\forall k=1,\dots,n$}
  \RightLabel{\textbf{(Coind-I)}}
  \TrinaryInfC{$\Delta\vdash$ corec
    $\vv{(\Gamma_k,y_k).g_k}:(\Gamma,y:C@id_\Gamma)\rat \nu@id_\Gamma$}
  \DisplayProof
\end{center}
gets rewritten to:
\begin{scprooftree}{0.95}
  \AxiomC{$\vdash C:\Gamma\rat*$}
  \AxiomC{\graybox{$\vdash\Gamma\equiv_\beta \Gamma'[\vv{D}]$}}
  \noLine
  \UnaryInfC{\graybox{$\vv{\vdash B_k\equiv_\beta A_k[\vv{D}][C/X]}$}}
  \AxiomC{\graybox{$\Phi\mid\Theta\mid\Delta\vdash \vv{D}:\Phi'$}}
  \noLine
  \UnaryInfC{$\vv{$\graybox{$\Phi\mid\mid$}$\Delta,\Gamma_k$\graybox{$[\vv{D}]$}$,y_k:(C@\sigma_k$\graybox{$[\vv{D}]$}$)\vdash g_k:\text{\graybox{$B_k$}}}$}
  \RightLabel{\textbf{(Coind-I)}}
  \TrinaryInfC{\graybox{$\Phi\mid\Theta\mid$}$\Delta\vdash$ corec \graybox{C to $\nu\langle\vv{D}\rangle$};
    $\vv{\text{Destr}_k\vv{x_k}\text{ } y_k = g_k}:(\Gamma,y:C@id_\Gamma)\rat \nu$\graybox{$[\vv{D}]$}$@id_\Gamma$}
 \end{scprooftree}

A corecursion produces a coinductive type \(\nu\).  We have to give it a type
\(C\) and list the destructors together with the expression they should be
destructed to. We get the syntax-directed rule analog as in the case of
recursion.

\section{Evaluation}
\label{sec:orgf9c7a7b}
There are two kinds of reduction steps in this system. The implementation of
this is in $\mathtt{Eval.hs}$. Will give the formal definition in the following.

The first is a reduction on the type level (written $\mathtt{\longrightarrow}$). It is defined
as follows:
\begin{align*}
  ((x).A) @ t \longrightarrow_p A[t/x]\\
\end{align*}
It is standard beta reduction. If we apply a lambda \((x).A)\) to a term \(t\) we
substitute this term for the binding variable \(x\) in the body. This body is
then the result of the reduction.

The other is the reduction on the term level (written $\mathtt{\succ}$). To define this
reduction, we need a action on types (written \(\widehat{C}(A)\)) and terms
(written \(\widehat{C}(t)\)), where the following holds.
\begin{prooftree}
  \AxiomC{$X : \Gamma_1\rat*\mid\Gamma_2'\vdash C:\Gamma_2\rat*$}
  \AxiomC{$\Gamma_1, x:A\vdash t:B$}
  \BinaryInfC{$\Gamma_2',\Gamma_2,x:\widehat{C}(A)\vdash\widehat{C}(t):\widehat{C}(B)$}
\end{prooftree}
Here, we have a type \(C\) with a free type variable \(X\) and a term \(t\) of type
\(B\) with a free term variable \(x\) of type \(A\). If we use the action of this
type on \(t\) we get a term with a type of this action on \(B\). This term
contains a free term variable \(x\) of type, the action applied to \(A\). The
type action is implemented in the module $\mathtt{TypeAction.hs}$. Both the type
action and the evaluation are done in the $\mathtt{Eval}$\;monad. This monad has access
to the previously defined declarations. We will now define the type action.

\begin{definition}
  \text{Let} $n \in \mathbb{N}$ \text{and} $1 \leq i \leq n$.
  \text{Let:}
  \begin{align*}
    X_1 : \Gamma_1 \rat \ast,\ldots,X_n : \Gamma_n \rat \ast\ \mid\ \Gamma' \vdash C : \Gamma \rat \ast \\
    \Gamma_i \vdash A_i : \ast \\
    \Gamma_i \vdash B_i : \ast \\
    \Gamma_i, x : A_i \vdash t_i : B_i
  \end{align*}
  \text{Then, we define the type action on terms inductively over $C$.}
  \begin{align*}
    \begin{array}{ll}
      \widehat{C}(\vv{t},t_{n+1}) = \widehat{C}(\vv{t})
      &\text{for \textbf{(TyVarWeak)}}\\
      \widehat{X_i}(\vv{t})=t_i\\
      \widehat{C'@s}(\vv{t})=\widehat{C'}(\vv{t})[s/y],
      &\text{for }\Theta\mid\Gamma'\vdash C':(y,\Gamma)\rat*\\
      \widehat{(y).C'}(\vv{t})=\widehat{C'}(\vv{t}),
      &\text{for }\Theta\mid(\Gamma',y)\vdash C':\Gamma\rat*\\
      \widehat{\mu(Y:\Gamma\rat*;\vv{\sigma};\vv{D})} =\text{rec}^{R_A}\vv{(\Delta_k,x).g_k}@\id{\Gamma}@x
      &\text{for } \Theta,Y:\Gamma\rat*\mid\Delta_k\vdash D_k:*\\
      \quad\text{with } g_k = \alpha_k^{R_B}@\id{\Delta_k}@\left(\widehat{D_k}(\vv{t},x)\right)\\
      \quad\text{and } R_A=\mu(Y:\Gamma\rat*;\vv{\sigma};\vv{D}[\vv{(\Gamma_i).A}/\vv{X}])\\
      \quad\text{and } R_B=\mu(Y:\Gamma\rat*;\vv{\sigma};\vv{D}[\vv{(\Gamma_i).B}/\vv{X}])\\
      \widehat{\nu(Y:\Gamma\rat*;\vv{\sigma};\vv{D})} =\text{corec}^{R_B}\vv{(\Delta_k,x).g_k}@\id{\Gamma}@x
      &\text{for } \Theta,Y:\Gamma\rat*\mid\Delta_k\vdash D_k:*\\
      \quad\text{with } g_k = \widehat{D_k}(\vv{t},x)[(\xi_k^{R_A}@\id{\Delta_k}@x)/x]\\
      \quad\text{and } R_A=\mu(Y:\Gamma\rat*;\vv{\sigma};\vv{D}[\vv{(\Gamma_i).A}/\vv{X}])\\
      \quad\text{and } R_B=\mu(Y:\Gamma\rat*;\vv{\sigma};\vv{D}[\vv{(\Gamma_i).B}/\vv{X}])\\
    \end{array}
  \end{align*}
  \text{And the type action on types as follows:}
  \begin{equation*}
    \widehat{C}(\vv{A})=C[\vv{(\Gamma_i).A}/\vv{X}]@\id{\Gamma}
  \end{equation*}
\end{definition}

The type action generates a term with a free variable \(x\).  In the type of
this term, we have changed all the free variables to the types of \(\vv{t}\).
We will show the proof in appendix \ref{sec:org421f57e}.

The reduction on terms  is subdivided into a reduction on recursion and one
on corecursion. Here, \(\sigma_k \bullet\tau\) is a context morphism, where we
first substitute with \(\tau\) and then with \(\sigma_k\).

The reduction on recursion is defined as follows:
\begin{align*}
  \rec \vv{(\Gamma_k,y_k).g_k}@(\sigma_k\bullet\tau)@(\alpha_k@\tau@u)\succ g_k\left[ \widehat{A_k}(\rec\vv{(\Gamma_k,y_k).g_k}@\id{\Gamma}@x)/y_k \right][\tau,u]\\
\end{align*}
If we apply a recursion \(\rec \vv{(\Gamma_k,y_k).g_k}\) to this context
morphism and a constructor \(\alpha_k@\tau@u\), which is fully applied, we
lookup the case for this constructor. In this case, we substitute \(\tau\) for
the variables from \(\Gamma_k\) and \(u\), where we apply the recursion to all
recursive occurrences, for \(y_k\). For this application, we need the type
action. So a recursion is destructing an inductive type and all its recursive
occurrences to another type, while we use different cases for the different
constructors of the type.

On the contrary, corecursion is constructing a coinductive type. It is defined
as follows:
\begin{align*}
  \xi_k@\tau@(\corec \vv{(\Gamma_k,y_k).g_k}@(\sigma_k\bullet\tau)@u)\succ \widehat{A_k}(\corec\vv{(\Gamma_k,y_k).g_k}@\id{\Gamma}@x)[g_k/x] [\tau,u]
\end{align*}
If we apply a destructor together with its arguments for its context
,\(\xi_k@\tau\), on such a construction \((\corec
   \vv{(\Gamma_k,y_k).g_k}@(\sigma_k\bullet\tau)@u)\), we are taking the case of
this destructor. In this case, we are applying the corecursion to all
recursive occurrences. \(\tau\) and \(u\) are substituted as in recursion.

\chapter{Examples}
\label{sec:orgc5acffd}

In this Section, we reiterate the example types from the paper.  We use our
syntax, which is defined in \ref{sec:org9120160}.  We will also show some functions
on these types.  On some of them, we will show the reduction steps in detail.

\section{Terminal and Initial Object}
\label{sec:org32f8072}

The terminal object is a type that has exactly one value. In category
theory, every object in the category has a unique morphism to it. We define
it as a coinductive type $\mathtt{Terminal}$\;with no destructors. It gets
a terminal and returns a terminal. To get a terminal value we use
corecursion on the unit type, which is the first-class terminal object.
\begin{lstlisting}
codata Terminal : Set where
terminal = corec Unit to Terminal where @ $\lozenge$
\end{lstlisting}
Contrary to the definition in the paper there is no destructor $\mathtt{Terminal}$.
In the paper definitions of coinductive or inductive types need at least one
de/constructor.  Therefore, our definition wouldn't work.

The initial object is a type that has no values. In category theory it is
the object which has a unique morphism to every other object in the category.
We define it inductively as $\mathtt{Intial}$\;with no constructor. In the paper, it is
defined with one constructor. This constructor want's one value of the same
type. We can't have a value of this type, because to get one we already need
one. Our way of defining it is shorter and more clear. We can't construct a
value of this type because we have no constructors. If we could get something
of type $\mathtt{Intial}$, we could generate with $\mathtt{exfalsum}$\;a value of arbitrary type
$\mathtt{C}$.
\begin{lstlisting}
data Initial : Set where
exfalsum$\langle$C : Set$\rangle$ = rec Initial to C where
\end{lstlisting}



\section{Natural Numbers and Extended Naturals}
\label{sec:orgc62542e}

We use the approach of Peano to define natural numbers.  Therefore, we use
the inductive type $\mathtt{Nat}$\;with the constructors $\mathtt{Zero}$\;and $\mathtt{Suc}$. $\mathtt{Zero}$\;is
just the number zero. Every constructor has to have an argument, which can
contain a recursive occurrence. Every Type $\mathtt{A}$\;is isomorphic to the
function type $\mathtt{Terminal\;\rat\;A}$. So we use $\mathtt{Terminal}$\;for this occurrence.
$\mathtt{Suc}$\;is the successor. So the meaning of $\mathtt{Suc\;n}$\;is \(n+1\).
\begin{lstlisting}
data Nat : Set where
   Zero : Terminal $\rat$ Nat
   Suc : Nat $\rat$ Nat
zero = Zero @ $\lozenge$
one = Suc @ zero
\end{lstlisting}
We can then define an identity recursion on it to see how reduction works.
It's a recursion that goes from $\mathtt{Nat}$\;to $\mathtt{Nat}$\;and gives back in every
case its input.
\begin{lstlisting}
id = rec Nat to Nat where
       Zero u = Zero @ u
       Succ n = Succ @ n
\end{lstlisting}

We use it on one to see all cases.
\begin{lstlisting}
id @ one = id @ (Succ @ zero)
         $\succ$ Succ @ n[$\widehat{X}$(id @ x)/n] [zero]
         = Succ @ $\widehat{X}$(id @ x) [zero]
         = Succ @ (id @ x)[zero]
         = Succ @ (id @ zero)
         = Succ @ (id @ (Zero @ $\lozenge$))
         $\succ$ Succ @ (Zero @ u[$\widehat{\text{Unit}}$(id @ x)/u][$\lozenge$])
         = Succ @ (Zero @ u[$\widehat{\text{Unit}}$(id @ x)/u][$\lozenge$])
         = Succ @ (Zero @ $\widehat{\text{Unit}}$(id @ x)[$\lozenge$])
         = Succ @ (Zero @ x)[$\lozenge$]
         = Succ @ (Zero @ x) = Succ @ zero = one
\end{lstlisting}
As expected the identity recursion applied to one gives back one.

We will now define extended naturals. There are also called co-natural numbers.
There are natural numbers with an additional value, infinity. We define it
coinductively with the predecessor as its only destructor. The predecessor is
either not defined or another natural number. We use the type $\mathtt{Maybe}$\;to
describe something which is either present (the constructor $\mathtt{Just}$) or
absent(the constructor $\mathtt{Nothing}$). We can define the successor as a
corecursion. The predecessor of the successor of $\mathtt{x}$\;is just $\mathtt{x}$. So the only
case of $\mathtt{corec}$\;returns a $\mathtt{Just\;x}$\;(remember $\mathtt{Prec}$\;returns a $\mathtt{Maybe\langle\;Conat\rangle}$
not a $\mathtt{Conat}$).
\begin{lstlisting}
data Maybe$\langle$A : Set$\rangle$ : Set where
  Nothing : Unit $\rat$ Maybe
  Just : A $\rat$ Maybe
nothing$\langle$A$\rangle$ = Nothing$\rangle$A$\rangle$ @ $\lozenge$
codata Conat : Set where
  Prec : Conat $\rat$ Maybe$\langle$Conat$\rangle$
succ = corec Conat to Conat where
         Prec x = Just$\langle$Conat$\rangle$ @ x
\end{lstlisting}
We now define the values zero and infinity.

\begin{lstlisting}
zero = (corec Unit to Conat where
          {Prev x = nothing$\langle$Unit$\rangle$}) @ $\lozenge$
infinity = (corec Unit to Conat where
              {Prev x = Just$\langle$Conat$\rangle$ @ x}) @ $\lozenge$
\end{lstlisting}

For $\mathtt{zero}$\;the predecessor is absent, there is no predecessor of 0 in the
natural numbers, so we give pack $\mathtt{Nothing}$.  We then have to apply the
$\mathtt{corec}$\;to $\mathtt{\lozenge}$\;to get the value.  The predecessor of $\mathtt{infinity}$\;should also
be $\mathtt{infinity}$.  We apply the $\mathtt{corec}$\;to another $\mathtt{Conat}$, so the $\mathtt{x}$\;is also a
$\mathtt{Conat}$.  We will now see that the predecessor on these values gives back the
right value.

{\scriptsize
\begin{flalign*}
  \Prev \apply \zero &\succ \widehat{\Maybe\langle X\rangle} \left(\underbrace{
     \begin{subarray}{l}
       \corecT{\Unit}{\Conat}\\
       \quad\{\;\Prev\; x = \nothing\langle\Unit\rangle\;\}
     \end{subarray}}_{t_1}\apply x\right) [\nothing\langle\Unit\rangle/x][\lozenge]&\\
              &= \recT{\Maybe\langle\Unit\rangle}{\Maybe\langle\Conat\rangle}&\\
              &\quad\quad \{ \Nothing\; u = \Nothing\langle\Conat\rangle \apply \widehat{\Unit}(t_1,u)&\\
              &\quad\quad  \;\Just\; c = \Just\langle\Conat\rangle \apply \widehat{X}(t_1,c)\} \apply x [\nothing\langle\Unit\rangle/x][\lozenge]&\\
              &= \underbrace{\begin{subarray}{l}
                \recT{\Maybe\langle\Unit\rangle}{\Maybe\langle\Conat\rangle} \\
                \quad \{\; \Nothing\; u = \Nothing\langle\Conat\rangle \apply u \\
                \quad \;\;\Just\; c = \Just\langle\Conat\rangle \apply t_1 \}
               \end{subarray}}_{t_2} \apply \nothing\langle\Unit\rangle&\\
              &\succ \Nothing\langle\Conat\rangle \apply u[\widehat{\Unit}(t_2 \apply x)/u][\lozenge]&\\
              &= \Nothing\langle\Conat\rangle \apply u[x/u][\lozenge]&\\
              &= \Nothing\langle\Conat\rangle \apply \lozenge&
\end{flalign*}
\begin{flalign*}
  \Prev \apply \infinity &\succ \widehat{\Maybe\langle X\rangle} \left(\underbrace{
     \begin{subarray}{l}
       \corecT{\Unit}{\Conat}\\
       \quad\{\; \Prev\; x = \Just\langle\Unit\rangle \apply x\; \}
     \end{subarray}}_{t_1} \apply x \right) [\Just\langle\Unit\rangle\apply/x][\lozenge]&\\
                  &= \recT{\Maybe\langle\Unit\rangle}{\Maybe\langle\Conat\rangle}&\\
                  &\quad\quad \{\; \Nothing\; u = \Nothing\langle\Conat\rangle \apply \widehat{\Unit}(t_1,u)&\\
                  &\quad\quad \;\;\Just\; c = \Just\langle\Conat\rangle \apply \widehat{X}(t_1,c)\;\} \apply x [\Just\langle\Unit\rangle\apply/x][\lozenge]&\\
                  &= \underbrace{\begin{subarray}{l}
                    \recT{\Maybe\langle\Unit\rangle}{\Maybe\langle\Conat\rangle} \\
                    \quad \{\; \Nothing\; u = \Nothing\langle\Conat\rangle \apply u \\
                    \quad \;\;\Just\; c = \Just\langle\Conat\rangle \apply t_1\; \}
                   \end{subarray}}_{t_2} \apply \Just\langle\Unit\rangle\apply&\\
                  &\succ \Just\langle\Conat\rangle \apply t_1[\widehat{\Unit}(t_2 \apply x)/x][\lozenge]&\\
                  &= \Just\langle\Conat\rangle \apply t_1[x/x][\lozenge]&\\
                  &= \Just\langle\Conat\rangle \apply \infinity&
\end{flalign*}}

\section{Binary Product and Coproduct}
\label{sec:org9ca13e8}

The product is defined as a coinductive type.  It has two destructors. The
first gives back the first element. And the second the second. To use this
type, the types A and B have to be instantiated to concrete types. We don't
have type polymorphism in our language. We also define a pair expression
which generates a pair over corecursion.
\begin{lstlisting}
codata Product$\langle$A : Set, B : Set$\rangle$ : Set where
   Fst : Product $\rat$ A
   Snd : Product $\rat$ B
pair$\langle$A : Set, B : Set$\rangle$ (x:A, y:B) = corec Unit where
                                      { Fst u $\rat$ x
                                      ; Snd u $\rat$ y} @ $\lozenge$
\end{lstlisting}
For types with other contexts, we have to define different product types.  For
example, if $\mathtt{B}$\;depends on $\mathtt{Nat}$, we define the product like the following:
\begin{lstlisting}
codata Pair$\langle$A : Set, B : (n : Nat) $\rat$ Set$\rangle$ : (n : Nat) $\rat$ Set where
  First : (n : Nat) $\rat$ Pair n $\rat$ A
  Second : (n : Nat) $\rat$ Pair n $\rat$ B @ n
\end{lstlisting}
Here, the product also depends on $\mathtt{Nat}$. If $\mathtt{A}$\;or $\mathtt{B}$\;depends on values the
product must also depend on these values. This is the product, which is used
for the definition of vectors in \cite{basold2016type}.

On $\mathtt{Product}$\;we can define the swap function.
\begin{lstlisting}
swap$\langle$A : Set, B : Set$\rangle$ =
  corec Product$\langle$A,B$\rangle$ to Product$\langle$B,A$\rangle$ where
         Fst x $\rat$ Snd x
         Snd x $\rat$ Fst x
\end{lstlisting}
This is a well-typed function as shown by the following proof
\begin{scprooftree}{0.95}
\AxiomC{$(A : *, B : *)\mid\mid\vdash$ Product$\langle$A,B$\rangle$ : $*$}
\AxiomC{$(A : *, B : *)\mid\mid(x:A) \vdash$ Snd @ x : Product$\langle$A,B$\rangle$ \circled{a}}
 \noLine
 \UnaryInfC{$(A : *, B : *)\mid\mid(y : B) \vdash$ Fst @ y : Product$\langle$A,B$\rangle$ \circled{b}}
 \BinaryInfC{$(A : *, B : *)\mid\mid \vdash$ swap : (p : Product$\langle$A,B$\rangle$) $\rat$ Product$\langle$B,A$\rangle$}
 \end{scprooftree}
We show \circled{a} in the following proof.  \circled{b} works analog.
\begin{prooftree}
\AxiomC{$(A : *, B : *)\mid\mid(x : A) \vdash$ Snd : $(x : A) \rat$ Product$\langle$A,B$\rangle$}
\AxiomC{$(x : A )(x)\rightsquigarrow A$}
\UnaryInfC{$(x : A) \vdash x : A$}
\BinaryInfC{$(A : *, B : *)\mid\mid(x : A) \vdash$ Snd @ x : Product$\langle$A,B$\rangle$}
\end{prooftree}
For brevity, we omitted the beta equality premises and the checking for of
the parameters. The beta equality premises wouldn't be interesting because
they all already syntactically identical.

The Binary Coproduct corresponds to the $\mathtt{Either}$\;type in Haskell.  It is defined
as an inductive type.  It is either $\mathtt{A}$\;or $\mathtt{B}$.  We have one constructor $\mathtt{Left}$
for $\mathtt{A}$\;and one constructor Right for $\mathtt{B}$.
\begin{lstlisting}
data Coproduct$\langle$A,B$\rangle$ : Set where
   Left : A $\rat$ Coproduct
   Right : B $\rat$ Coproduct
\end{lstlisting}

\section{Sigma and Pi Type}
\label{sec:orgff15342}

The sigma type is a dependent pair of two types.  The second type can depend on
the value of the first type.  It corresponds to exists in logic.  We define
it as an inductive type and call the constructor $\mathtt{Exists}$.
\begin{lstlisting}
data Sigma$\langle$A : Set,B : (x : A) $\rat$ Set$\rangle$ : Set where
   Exists : (x:A) $\rat$ B x $\rat$ Sigma
\end{lstlisting}

The pi type is a generalization of the function type to dependent types.
The type of the codomain or result of a function can depend on the value
We define it as a coinductive type.  To destruct a function we just apply
it to a value.  So the destructor is $\mathtt{Apply}$.

\begin{lstlisting}
codata Pi$\langle$A : Set,B : (x : A) $\rat$ Set$\rangle$ : Set where
   Apply : (x : A) $\rat$ Pi x $\rat$ B
\end{lstlisting}

To construct a function we use corecursion on $\mathtt{Unit}$. The identity function
is defined like this
\begin{lstlisting}
id$\langle$A : Set$\rangle$ = corec Unit to Pi$\langle$A,(v:A).A$\rangle$ where
       { Apply v p = v } @ $\lozenge$
\end{lstlisting}

Evaluation on one goes as follows:

\begin{lstlisting}
apply = Apply$\langle$Nat,(v : Nat).Nat$\rangle$
one = S @ (Z @ )
apply @ id$\langle$Nat$\rangle$ @ one
= apply @ one @ ((corec Unit to Pi$\langle$Nat,(x:Nat).Nat$\rangle$ where
                    Apply v p = v ) @ $\lozenge$)
$\succ \widehat{\text{Nat}} \left(\underbrace{
   \begin{subarray}{c}
     \text{corec Unit to Pi where} \\
     \text{  \{Apply' v \_ = v\} @ x}
   \end{subarray}}_t\right)$[v/x][one,$\lozenge$]
= (rec Nat to Nat where
     Zero x = Zero @ ($\widehat{\text{Unit}}$(t,x))
     Succ x = Suc @ ($\widehat{Y}$(t,x)))@x[v/x][one,$\lozenge$]
= (rec Nat to Nat where
     Zero x = Zero @ ($\widehat{\text{Unit}}$(t))
     Succ x = Suc @ x)@x[v/x][one,$\lozenge$]
= (rec Nat to Nat where
     Zero x = Zero @ ($\widehat{\text{Unit}}$())
     Succ x = Suc @ x) @ x[v/x][one,$\lozenge$]
= (rec Nat to Nat where
     Zero x = Zero @ x
     Succ x = Suc @ x) @ x[v/x][one,$\lozenge$]
= (rec Nat to Nat where
     Zero x = Zero @ x
     Succ x = Suc @ x) @ v[one,$\lozenge$]
= (rec Nat to Nat where
     Zero x = Zero @ x
     Succ x = Suc @ x) @ one
= one
\end{lstlisting}

\section{Vectors and Streams}
\label{sec:orgee29c18}

Vectors are a standard example for dependent types. They are like lists,
except their type depends on their length. For example, a vector $\mathtt{[1;2]}$\;has
type $\mathtt{Vector\langle\;Nat\rangle\;2}$, because its length is 2. It has 2
constructors $\mathtt{Nil}$\;and $\mathtt{Cons}$\;like lists. $\mathtt{Nil}$\;gives back the empty vector.
Because the length of the empty vector is zero its return type is $\mathtt{Vector\;0}$.
The second constructor $\mathtt{Cons}$\;takes a natural number $\mathtt{k}$, a value of type $\mathtt{A}$
and a vector of length $\mathtt{k}$, a $\mathtt{Vector\;k}$. It returns a new vector. Its head
is the first argument and its tail the second. So the length of the result is one
more than the second argument. Therefore, it is $\mathtt{Vector\;(Suc\;k)}$.  In
\cite{basold2016type} the head and tail are encoded in a pair.

\begin{lstlisting}
data Vector$\langle$A : Set$\rangle$ : (n:Nat) $\rat$ Set where
  Nil : Unit $\rat$ Vector zero
  Cons : (k:Nat, v:A) $\rat$ Vector @ k $\rat$ Vector (Suc @ k)
nil$\langle$A : Set$\rangle$ = Nil$\langle$A : Set$\rangle$ @ $\lozenge$
\end{lstlisting}

The function $\mathtt{extend}$\;takes a value $\mathtt{x}$\;and extends it to a vector.
\begin{lstlisting}
extend$\langle$A : Set$\rangle$ =
  rec Vec$\langle$A$\rangle$ to ((x).Vec$\langle$ A$\rangle$ @ (Suc x) where
    Nil u = Cons$\langle$A$\rangle$ @ x @ nil$\langle$A$\rangle$
    Cons k v = Cons$\langle$A$\rangle$ @ (Suc @ k) @ v
\end{lstlisting}
The type checking of this function goes as follows:
\begin{scprooftree}{0.75}
\AxiomC{(A : Set)$\mid\mid\vdash$ (x).(Vec$\langle$A$\rangle$ @ (Suc @ x)) : (k: Nat) $\rat$ *}
\noLine
\UnaryInfC{(A: Set)$\mid\mid$(u : A) $\vdash$ Cons$\langle$A$\rangle$ @ 0 @ (Nil$\langle$A$\langle$ @ ) : (x).(Vec$\langle$A$\rangle$ @ (Suc @ x)) @ 0}
\noLine
\UnaryInfC{(k : Nat, v : (x).(Vec @ (Suc @ x)) @ k) $\vdash$ Cons$\rangle$A$\langle$ @ (Suc @ k) @ v : (x).(Vec @ (Suc @ x)) @ (Suc @ k)}
\UnaryInfC{$\vdash$ extend$\langle$A$\rangle$ : (k:Nat,y : Vec$\langle$A$\rangle$ @ k) $\rat$ (x).(Vec$\langle$A$\rangle$ @ (Suc x)) @ k}
\end{scprooftree}
As an example, we evaluate a vector of length 1 with this function.  We choose length one
to see all $\mathtt{rec}$\;cases.
{\scriptsize
\begin{flalign*}
  &\extend\langle\Nat\rangle \apply 1 \apply (\Cons\langle\Nat\rangle \apply 0 \apply 0 \apply \nil\langle\Nat\rangle)&\\
  &= \extend\langle\Nat\rangle\apply(\Suc \apply k \bullet 0) \apply (\Cons\langle\Nat\rangle \apply 0 \apply 0 \apply \nil\langle\Nat\rangle)&\\
  &\succ \Cons\langle\Nat\rangle \apply (\Suc \apply k) \apply v \left[ \widehat{X\apply k}(\extend\langle\Nat\rangle \apply n \apply x)/v \right][0,\nil\langle\Nat\rangle]&\\
  &= \Cons\langle\Nat\rangle \apply (\Suc \apply k) \apply v \left[ \widehat{X}(\extend\langle\Nat\rangle \apply n \apply x)[k/n]/v \right][0,\nil\langle\Nat\rangle]&\\
  &= \Cons\langle\Nat\rangle \apply (\Suc \apply k) \apply v \left[ \extend \apply n \apply x[k/n]/v \right][0, \nil\langle\Nat\rangle]&\\
  &= \Cons\langle\Nat\rangle \apply (\Suc \apply k) \apply v \left[ \extend \apply k \apply x/v \right][0, \nil\langle\Nat\rangle]&\\
  &= \Cons\langle\Nat\rangle \apply (\Suc \apply k) \apply (\extend \apply k \apply x) [0,\nil\langle\Nat\rangle]&\\
  &= \Cons\langle\Nat\rangle \apply (\Suc \apply 0) \apply (\extend \apply 0 \apply (\nil\langle\Nat\rangle))&\\
  &= \Cons\langle\Nat\rangle \apply 1 \apply (\extend \apply 0 \apply (\Nil\langle\Nat\rangle \apply ))&\\
  &\succ \Cons\langle\Nat\rangle \apply 1 \apply (\Cons\langle\Nat\rangle \apply 0 \apply (\Nil\langle\Nat\rangle \apply ))\left[ \widehat{\Unit}(\extend \apply k \apply x) / u  \right][\lozenge]&\\
  &= \Cons\langle\Nat\rangle \apply 1 \apply (\Cons\langle\Nat\rangle \apply 0 \apply (\Nil\langle\Nat\rangle \apply x))[\lozenge]&\\
  &= \Cons\langle\Nat\rangle \apply 1 \apply (\Cons\langle\Nat\rangle \apply 0 \apply (\Nil\langle\Nat\rangle \apply ))&
\end{flalign*}}
Here, we write \(1\) for $\mathtt{Suc\;@\;(Zero\;@\;)}$\;and \(0\) for $\mathtt{Zero\;@\;\lozenge}$.

With the help of extended naturals, we can define partial streams. Those are
streams that depend on their definition depth. Like non-dependent streams,
they are coinductive and have 2 destructors for head and tail.
\begin{lstlisting}
codata PStr$\langle$A : Set$\rangle$: (n: ExNat) $\rat$ Set where
   hd : (k : ExNat) $\rat$ PStr$\langle$A$\rangle$ (succE k) $\rat$ A
   tl : (k : ExNat) $\rat$ PStr$\langle$A$\rangle$ (succE k) $\rat$ PStr$\langle$A$\rangle$ @ k
\end{lstlisting}
These streams are like vectors except they also can be infinite long. This is
in contrary to non-dependent streams.  A non-dependent stream could not be of
length zero.  Because then a call of $\mathtt{hd}$\;and $\mathtt{tl}$\;on it wouldn't be defined.  In the
dependent case, the type checker wouldn't allow such a call because $\mathtt{hd}$\;and
$\mathtt{tl}$\;expect streams which are at least of length one.  We can then define
$\mathtt{repeat}$.
\begin{lstlisting}
repeat$\langle$A : Set$\rangle$(x : A, n : Conat) =
  corec (n : Conat).Unit to PStr$\langle$A$\rangle$ where
    { Hd k s = x
    ; Tl k s = $\lozenge$ } @ n @ $\lozenge$
\end{lstlisting}
This function gets a value and an extended natural number. It generates a
constant partial stream of that value with the number as its length.


\chapter{Conclusion}
\label{sec:org835b5ea}
We have implemented a dependent type theory with inductive and coinductive types.
In this theory, contrary to Coq and Agda, coinductive types can also depend on
values. Contrary to the theory of the paper we can define schemata like
$\mathtt{Maybe\langle\;A\;:\;Set\rangle}$\;\linebreak where $\mathtt{A}$\;can be an arbitrary type
of kind $\mathtt{Set}$.

One downside is that we don't have universes. This prevents type polymorphism.
Further work needs to be done to solve this. Another problem is, that each
constructor or destructor has at least one argument. The argument with the
recursive occurrence. For example, we have to apply a unit to the constructors
of a boolean type. We could allow recursive occurrences in the contexts of the
constructors and destructors. This makes it possible to remove the argument
with the recursive occurrence. Then we have to change the evaluation rules.

Our system allowed us to define the (depended) function type. Therefore, we
don't have it as a primitive expression. We are hopeful, that in the future we
get a more mainstream language, like Coq or Agda, where the dependet function
is definable. As already mentioned in the introduction this would lead to a
symmetrical language.

\appendix

\chapter{Type action proof}
\label{sec:org421f57e}

\begin{theorem}
\((\Gamma).A@\id{\Gamma}\leftrightarrow_T A\)
\label{abstrid}
\end{theorem}
\begin{proof}
We show this by induction on the length of \(\Gamma\)
\begin{itemize}
\item \(\Gamma=\epsilon\):
\begin{equation*}
   A \longleftrightarrow_T A
\end{equation*}
\item \(\Gamma=x:B,\Gamma'\):
\begin{equation*}
  (x:B,\Gamma').A@x@\id{\Gamma'}
  \longrightarrow_p(\Gamma').A@\id{\Gamma'}[x/x]
  = (\Gamma').A@\id{\Gamma'} \overset{IdH.}{\longleftrightarrow_T}A
\end{equation*}
\end{itemize}
\end{proof}
\begin{theorem}
The following rule holds
\begin{prooftree}
\AxiomC{$x:A\vdash t:B$}
\AxiomC{$A\longleftrightarrow_TA'$}
\BinaryInfC{$x:A'\vdash t:B$}
\end{prooftree}
\label{ctxconv}
\end{theorem}
\begin{proof}
We show this by induction on t
\end{proof}
\begin{theorem}
The typing rule (5) in the paper holds
\begin{prooftree}
  \AxiomC{$X:\Gamma_1\rat*\mid\Gamma'\vdash C:\Gamma\rat*$}
  \AxiomC{$\Gamma_1,x:A\vdash t:B$}
  \BinaryInfC{$\Gamma',\Gamma,x:\widehat{C}(A)\vdash\widehat{C}(t):\widehat{C}(B) $}
\end{prooftree}
\end{theorem}
\begin{proof}
First we will generalize the rule to
\begin{prooftree}
  \AxiomC{$X_1:\Gamma_1\rat*,\dots,X_n:\Gamma_n\rat*\mid\Gamma'\vdash C:\Gamma\rat*$}
  \AxiomC{$\Gamma_i,x:A_i\vdash t_i:B_i$}
  \BinaryInfC{$\Gamma',\Gamma,x:\widehat{C}(\vv{A})\vdash\widehat{C}(\vv{t}):\widehat{C}(\vv{B}) $}
\end{prooftree}
Then, we gonna show it by Induction on the derivation \(\mathcal{D}\) of \(C\)
\begin{itemize}
\item $\mathcal{D}$ =
  \AxiomC{}
  \topI{$\vdash\top:*$}
  \DisplayProof

Then, the type actions got calculated as follows
\begin{align*}
  &\widehat{\top}(\vv{A}) = \widehat{\top}() = \top\\
  &\widehat{\top}(\vv{t}) = \widehat{\top}() = x\\
  &\widehat{\top}(\vv{B}) = \widehat{\top}() = \top
\end{align*}
We than got the following prooftree
\begin{prooftree}
  \AxiomC{$\vdash\top:*$}
  \RightLabel{\textbf{(Proj)}}
  \UnaryInfC{$x:\top\vdash x:\top$}
\end{prooftree}
\item $\mathcal{D}$ =
  \Di{1}
  \UnaryInfC{$X_1:\Gamma_1\rat*,\dots,X_{n-1}:\Gamma_{n-1}$\TyCtx}
  \Di{2}
  \UnaryInfC{$\Gamma_n$\Ctx}
  \TyVarI{$X_1:\Gamma_1\rat*,\dots,X_n:\Gamma_n\rat*\mid\emptyset\vdash X_n:\Gamma_n\rat*$}
  \DisplayProof

Again we calculate the type actions
\begin{align*}
  &\widehat{X_n}(\vv{A}) = X_n[\vv{(\Gamma_i).A}/\vv{X}]@\id{\Gamma_n}= X_n[(\Gamma_n).A_n/X_n]@\id{\Gamma_n} = (\Gamma_n).A_n@\id{\Gamma_n}\\
  &\widehat{X_n}(\vv{t}) = t_n\\
  &\widehat{X_n}(\vv{B}) = X_n[\vv{(\Gamma_i).B}/\vv{X}]@\id{\Gamma_n}= X_n[(\Gamma_n).B_n/X_n]@\id{\Gamma_n} = (\Gamma_n).B_n@\id{\Gamma_n}\\
\end{align*}
We know from the first premise that \(\Gamma=\Gamma_n\) and \(\Gamma'=\emptyset\)

Here, we got the prooftree
\begin{scprooftree}{0.9}
\AxiomC{$\Gamma_n,x:A\vdash t:B$}
\AxiomC{}
\RightLabel{Thrm. \ref{abstrid}}
\UnaryInfC{$A\longleftrightarrow_T(\Gamma_n).A@\id{\Gamma_n}$}
\RightLabel{Thrm. \ref{ctxconv}}
\BinaryInfC{$\Gamma_n,x:(\Gamma_n).A@\id{\Gamma_n}\vdash t:B$}
\AxiomC{}
\RightLabel{Thrm. \ref{abstrid}}
\UnaryInfC{$B\longleftrightarrow_T(\Gamma_n).B@\id{\Gamma_n}$}
\RightLabel{Conv}
\BinaryInfC{$\Gamma_n,x:(\Gamma_n).A@\id{\Gamma_n}\vdash t_n:(\Gamma_n).B@\id{\Gamma_1}$}
\end{scprooftree}

\item $\mathcal{D}$ =
  \Di{1}
  \UnaryInfC{$X_1:\Gamma_1\rat*,\dots,X_n:\Gamma_n\mid\Gamma'\vdash C:\Gamma\rat*$}
  \Di{2}
  \UnaryInfC{$\Gamma_n$\Ctx}
  \TyVarWeak{$X_1:\Gamma_1\rat*,\dots,X_{n+1}:\Gamma_{n+1}\rat*\mid\Gamma'\vdash C:\Gamma\rat*$}
  \DisplayProof

Here, we got the prooftree
\begin{prooftree}
  \AxiomC{$X_1:\Gamma_1\rat*,\dots,X_{n+1}:\Gamma_{n+1}\rat*\mid\Gamma'\vdash C:\Gamma\rat*$}
  \RightLabel{(*)}
  \UnaryInfC{$X_1:\Gamma_1\rat*,\dots,X_n:\Gamma_n\rat*\mid\Gamma'\vdash C:\Gamma\rat*$}
  \AxiomC{$\Gamma_i,x:A_i\vdash t_i:B_i$}
  \RightLabel{IdH.}
  \BinaryInfC{$\Gamma',\Gamma,x:\underbrace{\widehat{C}(\vv{A})}_{\overset{(**)}{=}\widehat{C}(\vv{A},A_{n+1})}\vdash\underbrace{\widehat{C}(\vv{t})}_{\overset{(***)}{=}\widehat{C}(\vv{t},t_{n+1})}:\underbrace{\widehat{C}(\vv{B})}_{\overset{(**)}{=}\widehat{C}(\vv{B},B_{n+1})} $}
\end{prooftree}

(\texttt{*}) Here, we undo \textbf{(TyVar-Weak)}

(\texttt{**}) \(X_{n+1}\) doesn't occur free in C, otherwise \(\mathcal{D}_1\) wouldn't be possible

(\texttt{***}) Case for \textbf{(TyVar-Weak)} of type actions on terms

\item $\mathcal{D}$ =
  \begin{scprooftree}{0.8}
  \Di{1}
  \UnaryInfC{$X_1:\Gamma_1\rat*,\dots,X_n:\Gamma_n\mid\Gamma'\vdash C:\Gamma\rat*$}
  \Di{2}
  \UnaryInfC{$X_1:\Gamma_1\rat*,\dots,X_n:\Gamma_n\mid\Gamma'\vdash D:*$}
  \TyWeak{$X_1:\Gamma_1\rat*,\dots,X_n:\Gamma_n\rat*\mid\Gamma',y:D\vdash C:\Gamma\rat*$}
  \end{scprooftree}

Here, we got the prooftree
\begin{scprooftree}{0.6}
  \AxiomC{$X_1:\Gamma_1\rat*,\dots,X_n:\Gamma_n\rat*\mid\Gamma',y:D\vdash C:\Gamma\rat*$}
  \RightLabel{(*)}
  \UnaryInfC{$X_1:\Gamma_1\rat*,\dots,X_n:\Gamma_n\rat*\mid\Gamma'\vdash C:\Gamma\rat*$}
  \AxiomC{$\Gamma_i,x:A_i\vdash t_i:B_i$}
  \RightLabel{IdH.}
  \BinaryInfC{$\Gamma',\Gamma,x:\widehat{C}(\vv{A})\vdash\widehat{C}(\vv{t}):\widehat{C}(\vv{B})$}
  \AxiomC{$X_1:\Gamma_1\rat*,\dots,X_n:\Gamma_n\mid\Gamma'\vdash D:*$}
  \TermWeak{$\Gamma',\Gamma,x:\widehat{C}(\vv{A})y\vdash\widehat{C}(\vv{t}):\widehat{C}(\vv{B})$}
\end{scprooftree}

(\texttt{*}) Here, we undo \textbf{(Ty-Weak)}

\item $\mathcal{D}= $
\AxiomC{$X_1:\Gamma_1,\ldots,X_n:\Gamma_n\mid\Gamma'\vdash C':(y:D,\Gamma)\rat* $}
\AxiomC{$\Gamma'\vdash s: D$}
\TyInst{$X_1:\Gamma_1,\ldots,X_n:\Gamma_n\mid\Gamma'\vdash C'@s:\Gamma\rat* $}
\DisplayProof

Then, we got the following induction hypothesis
\begin{prooftree}
  \AxiomC{$X_1:\Gamma_1\rat*,\dots,X_n:\Gamma_n\rat*\mid\Gamma'\vdash C':(y:D,\Gamma)\rat*$}
  \AxiomC{$\Gamma_i,x:A_i\vdash t_i:B_i$}
  \BinaryInfC{$\Gamma',y:D,\Gamma,x:\widehat{C'}(\vv{A})\vdash\widehat{C'}(\vv{t}):\widehat{C'}(\vv{B}) $}
\end{prooftree}

Calculated type actions:
\begin{align*}
  &\widehat{C'@s}(\vv{A})=C'@s[\vv{(\Gamma_i).A}/\vv{X}]@\id{\Gamma}=C'[\vv{(\Gamma_i).A}/\vv{X}]@s@\id{\Gamma}
  =\widehat{C'}(\vv{A})[s/y]\\
  &\widehat{C'@s}(\vv{t})=\widehat{C'}(\vv{t})[s/y]\\
  &\widehat{C'@s}(\vv{B})=C'@s[\vv{(\Gamma_i).B}/\vv{X}]@\id{\Gamma}=C'[\vv{(\Gamma_i).B}/\vv{X}]@s@\id{\Gamma}
  =\widehat{C'}(\vv{B})[s/y]\\
\end{align*}

We then got the following prooftree
\begin{scprooftree}{0.95}
  \AxiomC{$X_1:\Gamma_1\rat*,\dots,X_n\rat*\mid\Gamma_2'\vdash C'@s:\Gamma_2[s/y]\rat*$}
  \RightLabel{(*)}
  \UnaryInfC{$X_1:\Gamma_1\rat*,\dots,X_n:\Gamma_n\rat*\mid\Gamma_2'\vdash C':(y:D,\Gamma_2)\rat*$}
  \AxiomC{$\Gamma_i,x:A_i\vdash t_i:B_i$}
  \RightLabel{IdH.}
  \BinaryInfC{$\Gamma_2',y:D,\Gamma_2,x:\widehat{C'}(\vv{A})\vdash\widehat{C'}(\vv{t}):\widehat{C'}(\vv{B}) $}
  \UnaryInfC{$\Gamma_2',\Gamma_2[s/y],x:\widehat{C'}(\vv{A})[s/y]\vdash\widehat{C'}(\vv{t})[s/y]:\widehat{C'}(\vv{B})[s/y] $}
\end{scprooftree}
(\texttt{*}) This is the reverse of \textbf{(Ty-Inst)}.

\item $\mathcal{D}= $
\AxiomC{$X_1:\Gamma_1,\ldots,X_n:\Gamma_n\mid\Gamma',y:D\vdash C':\Gamma\rat* $}
\ParamAbstr{$X_1:\Gamma_1,\ldots,X_n:\Gamma_n\mid\Gamma'\vdash (y).C':(y:D,\Gamma)\rat* $}
\DisplayProof

Calculated type actions:
\begin{align*}
  \widehat{(y).C'}(\vv{A})&=(y).C'[\vv{(\Gamma_i.A)}/\vv{X}]@\id{\Gamma}\\
                     &=(y).(C'[\vv{(\Gamma_i.A)}/\vv{X}])@y@\id{\Gamma}\\
                     &\longleftrightarrow_T(C'[\vv{(\Gamma_i.A)}/\vv{X}])@\id{\Gamma}\\
                     &=\widehat{C'}(\vv{A})\\
  \widehat{(y).C'}(\vv{t})&=\widehat{C'}(\vv{t})\\
  \widehat{(y).C'}(\vv{B})&=(y).C'[\vv{(\Gamma_i.B)}/\vv{X}]@\id{\Gamma}\\
                     &=(y).(C'[\vv{(\Gamma_i.B)}/\vv{X}])@y@\id{\Gamma}\\
                     &\longleftrightarrow_T(C'[\vv{(\Gamma_i.B)}/\vv{X}])@\id{\Gamma}\\
                     &=\widehat{C'}(\vv{B})\\
\end{align*}

The prooftree then becomes the following
\begin{scprooftree}{0.95}
  \AxiomC{$X_1:\Gamma_1\rat*,\dots,X_n:\Gamma_n\rat*\mid\Gamma'\vdash (y).C':(y:D,\Gamma)\rat*$}
  \RightLabel{(*)}
  \UnaryInfC{$X_1:\Gamma_1\rat*,\dots,X_n:\Gamma_n\rat*\mid y:D,\Gamma'\vdash C':\Gamma\rat*$}
  \AxiomC{$\Gamma_i,x:A_i\vdash t_i:B_i$}
  \RightLabel{IdH.}
  \BinaryInfC{$y:D,\Gamma',\Gamma,x:\widehat{C'}(\vv{A})\vdash\widehat{C'}(\vv{t}):\widehat{C'}(\vv{B})$}
\end{scprooftree}
(\texttt{*}) This is the reverse of \textbf{(Param-Abstr)}.

\item \(\mathcal{D}\) =
\begin{prooftree}
    \Di{1}
    \UnaryInfC{$\sigma_k:\Delta_k\triangleright\Gamma$}
    \Di{2}
    \UnaryInfC{$X_1:\Gamma_1\rat*,\dots,X_n\rat*,X:\Gamma\rat*\mid\Delta_k\vdash D_k:*$}
    \FPTy
    \BinaryInfC{$X_1:\Gamma_1\rat*,\dots,X_{n}\rat*\mid\emptyset\vdash\mu(Y:\Gamma\rat*;\vv{\sigma};\vv{D}):\Gamma\rat*$}
\end{prooftree}
From this we know \(\Gamma'=\emptyset\)

Calculated type actions:
\begin{align*}
  &\widehat{\mu(Y:\Gamma\rat*;\vv{\sigma};\vv{D})}(\vv{A})\\
  &=\mu(Y:\Gamma\rat*;\vv{\sigma};\vv{D})[\vv{(\Gamma_i).A}/\vv{X}]@\id{\Gamma}\\
  &=\mu(Y:\Gamma\rat*;\vv{\sigma};\vv{D}[\vv{(\Gamma_i).A}/\vv{X}])@\id{\Gamma}\\
  &\widehat{\mu(Y:\Gamma\rat*;\vv{\sigma};\vv{D})}(\vv{t})\\
  &=\text{rec}^{\mu(Y:\Gamma\rat*;\vv{\sigma};\vv{D}[\vv{(\Gamma_i).A}/\vv{X}])}\vv{(\Delta_k,x).\alpha_k@\id{\Delta_k}@\widehat{D_k}(\vv{t},x)}@\id{\Gamma}@x\\
  &\widehat{\mu(Y:\Gamma\rat*;\vv{\sigma};\vv{D})}(\vv{B})\\
  &=\mu(Y:\Gamma\rat*;\vv{\sigma};\vv{D})[\vv{(\Gamma_i).B}/\vv{X}]@\id{\Gamma}\\
  &=\mu(Y:\Gamma\rat*;\vv{\sigma};\vv{D}[\vv{(\Gamma_i).B}/\vv{X}])@\id{\Gamma}
\end{align*}

From the assumptions
\begin{align*}
&X_1:\Gamma_1\rat*,\dots,X_n:\Gamma_n\rat*\mid\emptyset\vdash \mu(Y:\Gamma\rat*;\vv{\sigma};\vv{D}):\Gamma\rat*\\
&\Gamma_i,x:A_i\vdash t_i:B_i
\end{align*}
We have to proof that in \textbf{Ctx}
\begin{equation*}
 \Gamma,x:\mu(Y:\Gamma\rat*;\vv{\sigma};\vv{D}[\vv{(\Gamma_i).A}/\vv{B}])@\id{\Gamma}
\end{equation*}
the expression
\begin{equation*}
 \text{rec}^{\mu(Y:\Gamma\rat*;\vv{\sigma};\vv{D}[\vv{(\Gamma_i).A}/\vv{X}])}\vv{(\Delta_k,y).\alpha_k@\id{\Delta_k}@\widehat{D_k}(t,y)}@\id{\Gamma}@x
\end{equation*}
has type
\begin{equation*}
\mu(Y:\Gamma\rat*;\vv{\sigma};\vv{D}[\vv{(\Gamma_i).B}/\vv{X}])@\id{\Gamma}
\end{equation*}
We can use the induction hypothesis
\begin{prooftree}
  \AxiomC{$X_1:\Gamma_1\rat*,\dots,X_n:\Gamma_n\rat*,Y:\Gamma_{n+1}\rat*\mid\Delta_k\vdash D_k:*$}
  \AxiomC{$\Gamma_i,x:A_i\vdash t_i:B_i$}
  \BinaryInfC{$\Delta_k,x:\widehat{D_k}(\vv{A},A_{n+1})\vdash\widehat{D_k}(\vv{t},y):\widehat{D_k}(\vv{B},B_{n+1}) $}
\end{prooftree}
See \ref{sec:prc} for a proof of it.

\item \(\mathcal{D}\) =
\begin{prooftree}
    \Di{1}
    \UnaryInfC{$\sigma_k:\Delta_k\triangleright\Gamma$}
    \Di{2}
    \UnaryInfC{$X_1:\Gamma_1\rat*,\dots,X_n\rat*,X:\Gamma\rat*\mid\Delta_k\vdash D_k:*$}
    \FPTy
    \BinaryInfC{$X_1:\Gamma_1\rat*,\dots,X_{n}\rat*\mid\emptyset\vdash\nu(Y:\Gamma\rat*;\vv{\sigma};\vv{D}):\Gamma\rat*$}
\end{prooftree}
From this we know \(\Gamma'=\emptyset\).

Calculated type actions:
\begin{align*}
  &\widehat{\nu(Y:\Gamma\rat*;\vv{\sigma};\vv{D})}(\vv{A})\\
  &=\nu(Y:\Gamma\rat*;\vv{\sigma};\vv{D})[\vv{(\Gamma_i).A}/\vv{X}]@\id{\Gamma}\\
  &=\nu(Y:\Gamma\rat*;\vv{\sigma};\vv{D}[\vv{(\Gamma_i).A}/\vv{X}])@\id{\Gamma}\\
  &\widehat{\nu(Y:\Gamma\rat*;\vv{\sigma};\vv{D})}(\vv{t})\\
  &=\text{corec}^{\nu(Y:\Gamma\rat*;\vv{\sigma};\vv{D}[(\vv{\Gamma_i).B}/\vv{X}])}\vv{(\Delta_k,x)\widehat{D_k}(\vv{t},x)[(\xi_k@\id{\Delta_k}@x)/x]}@\id{\Gamma}@x\\
  &\widehat{\nu(Y:\Gamma\rat*;\vv{\sigma};\vv{D})}(\vv{B})\\
  &=\nu(Y:\Gamma\rat*;\vv{\sigma};\vv{D})[\vv{(\Gamma_i).B}/\vv{X}]@\id{\Gamma}\\
  &=\nu(Y:\Gamma\rat*;\vv{\sigma};\vv{D}[\vv{(\Gamma_i).B}/\vv{X}])@\id{\Gamma}
\end{align*}

From the assumptions
\begin{align*}
&X_1:\Gamma_1\rat*,\dots,X_n:\Gamma_n\rat*\mid\emptyset\vdash \nu(Y:\Gamma\rat*;\vv{\sigma};\vv{D}):\Gamma\rat*\\
&\Gamma_i,x:A_i\vdash t_i:B_i
\end{align*}
We have to proof that in \textbf{Ctx}
\begin{equation*}
 \Gamma,x:\nu(Y:\Gamma\rat*;\vv{\sigma};\vv{D}[(\Gamma_1).A/X])@\id{\Gamma}
\end{equation*}
the expression
\begin{equation*}
 \text{corec}^{\nu(Y:\Gamma\rat*;\vv{\sigma};\vv{D}[(\vv{\Gamma_i).B}/\vv{X}])}\vv{(\Delta_k,x)\widehat{D_k}(\vv{t},x)[(\xi_k@\id{\Delta_k}@x)/x]}@\id{\Gamma}@x\\
\end{equation*}
has type
\begin{equation*}
\nu(Y:\Gamma\rat*;\vv{\sigma};\vv{D}[\vv{(\Gamma_i).B}/\vv{X}])@\id{\Gamma}
\end{equation*}
We can use the induction hypothesis
\begin{prooftree}
  \AxiomC{$X_1:\Gamma_1\rat*,\dots,X_n:\Gamma_n\rat*,Y:\Gamma_{n+1}\rat*\mid\Delta_k\vdash D_k:*$}
  \AxiomC{$\Gamma_i,y_k:A_i\vdash t_i:B_i$}
  \BinaryInfC{$\Delta_k,y_k:\widehat{D_k}(\vv{A},A_{n+1})\vdash\widehat{D_k}(\vv{t},y):\widehat{D_k}(\vv{B},B_{n+1}) $}
\end{prooftree}
See \ref{sec:prc} for this proof.
\end{itemize}
\end{proof}

  \begin{landscape}
  \section{Proofs for recursion and corecursion}
  \label{sec:prc}
\begin{prooftree}
   \AxiomC{$\Gamma_1\vdash\sigma:\Gamma_2$}
   \AxiomC{$\Gamma_3\vdash\tau:\Gamma_1$}
   \RightLabel{($*$)}
   \BinaryInfC{$\Gamma_3\vdash\sigma\circ\tau:\Gamma_2$}
\end{prooftree}
      \begin{prooftree}
      \D
      \UnaryInfC{$\Delta,\Gamma_k,y_k:A_k[C/X]\vdash g_k:C@\sigma_k$}
      \RightLabel{TyAct}
      \UnaryInfC{$\Delta,\Gamma_k,x:A_k[\mu/X]\vdash g_k[\widehat{A_k}(\rec^\mu\vv{(\Gamma_k,y_k).g_k}@\id{\Gamma}@x/y_k]$}
      \D
      \UnaryInfC{$\Delta\vdash\tau:\Gamma_k$}
      \D
      \UnaryInfC{$\Delta\vdash u:A_k[\mu/X]$}
      \TrinaryInfC{$\Delta\vdash g_k[\widehat{A_k}(\rec^\mu\vv{(\Gamma_k,y_k).g_k}@\id{\Gamma}@x)/y_k][\tau,u]:C@\sigma_k$}
    \end{prooftree}
  \begin{scprooftree}{0.7}
    \D
    \UnaryInfC{$\Delta,\Gamma_k,y_k:\Delta_k[C/X]\vdash g_k: C@\sigma_k$}
    \IndE
    \UnaryInfC{$\Delta\vdash\rec^\mu\vv{(\Gamma_k,y_k).g_k}:(\Gamma,x:\mu@\sigma_k)\rat C@\sigma_k$}
    \D
    \UnaryInfC{$\Gamma_k\vdash\sigma_k:\Gamma$}
    \D
    \UnaryInfC{$\Delta\vdash\tau:\Gamma_k$}
    \RightLabel{($*$)}
    \BinaryInfC{$\Delta\vdash\sigma_k\circ\tau:\Gamma$}
    \Inst{$\Delta\vdash(rec^\mu\vv{(\Gamma_k,y_k).g_k}@(\sigma_k\circ\tau)):(x:\mu@\sigma_k)\rat C@\sigma_k$}
    \AxiomC{}
    \IndI{$\Delta\vdash\alpha_k^\mu:(\Gamma_k,y:A_k[\mu/X])\rat\mu@\sigma_k$}
    \D
    \UnaryInfC{$\Delta\vdash\tau:\Gamma_k$}
    \Inst{$\Delta\vdash\alpha_k^\mu @\tau: (y:A_k[\mu/X])\rat\mu@\sigma_k$}
    \D
    \UnaryInfC{$\Delta\vdash u:A_k[\mu/X])$}
    \Inst{$\Delta\vdash\alpha_k^\mu @\tau@u:\mu@\sigma_k$}
    \Inst{$\Delta\vdash(\rec^\mu\vv{(\Gamma_k,y_k).g_k}@(\sigma_k\circ\tau))@(\alpha_k^\mu @\tau@u):C@\sigma_k$}
  \end{scprooftree}
  \begin{scprooftree}{0.48}
    \AxiomC{$\vdash\overbrace{\mu[\vv{(\Gamma_k).A_k}/\vv{X_k}]}^C:\Gamma\rat*$}
    \AxiomC{}
    \IndI{$\Gamma,x:\mu[\vv{(\Gamma_k).A_k}/\vv{X_k}]@\id{\Gamma},\Delta_k,y_k:A_k[\mu[\vv{(\Gamma_k).B_k}/\vv{X_k}]/X]\vdash\alpha_k^{\mu[\vv{(\Gamma_k).B_k}/\vv{X_k}]}:(\Delta_k,y_k:A_k[\vv{(\Gamma_k).B_k}/\vv{X_k}])\rat\mu[\vv{(\Gamma_k).B_k}/\vv{X_k}]@\sigma_k$}
    \AxiomC{$\Gamma,x:\mu[\vv{(\Gamma_k).A_k}/\vv{X_k}]@\id{\Gamma},\Delta_k,y_k:\overbrace{A_k[\mu[\vv{(\Gamma_k).B_k}/\vv{X_k}]/X]}^{=\widehat{A_k}(\vv{B},?)}\vdash\widehat{A_k}(\vv{t},y_k):\overbrace{A_k[\mu[\vv{(\Gamma_k).B_k}/\vv{X_k}]/X]}^{=\widehat{A_k}(\vv{B},?)}$}
    \Inst{$\Gamma,x:\mu[\vv{(\Gamma_k).A_k}/\vv{X_k}]@\id{\Gamma},\Delta_k,y_k:A_k[\mu[\vv{(\Gamma_k).B_k}/\vv{X_k}]/X]\vdash\alpha_k^{\mu[\vv{(\Gamma_k).B_k}/\vv{X_k}]}@\id{\Delta_k}@\widehat{A_k}(\vv{t},y_k):\mu[\vv{(\Gamma_k).B_k}/\vv{X_k}]@\sigma_k$}
    \IndE
    \BinaryInfC{$\Gamma,x:\mu[\vv{(\Gamma_k).A_k}/\vv{X_k}]@\id{\Gamma}\vdash\rec^{\mu[\vv{(\Gamma_k).A_k}/\vv{X_k}]}\vv{(\Delta_k,y_k).\alpha_k^{\mu[\vv{(\Gamma_k).B_k}/\vv{X_k}]}@\id{\Delta_k}@\widehat{A_k}(\vv{t},y_k)}:(\Gamma,x:\mu[\vv{(\Gamma_k).A_k}/\vv{X_k}]@\id{\Gamma})\rat\mu[\vv{(\Gamma_k).B_k}/\vv{X_k}]@\id{\Gamma}$}
    \AxiomC{$\ldots$}
    \Inst{$\Gamma,x:\mu[\vv{(\Gamma_k).A_k}/\vv{X_k}]@\id{\Gamma}\vdash\rec^{\mu[\vv{(\Gamma_k).A_k}/\vv{X_k}]}\vv{(\Delta_k,y_k).\alpha_k^{\mu[\vv{(\Gamma_k).B_k}/\vv{X_k}]}@\id{\Delta_k}@\widehat{A_k}(\vv{t},y_k)}@\id{\Gamma}@x:\mu[\vv{(\Gamma_k).B_k}/\vv{X_k}]@\id{\Gamma}$}
  \end{scprooftree}
  \begin{scprooftree}{0.78}
    \AxiomC{$\vdash\overbrace{\nu[\vv{(\Gamma_k).A_k}/\vv{X_k}]}^C:\Gamma\rat*$}
    \AxiomC{$\Gamma,x:\nu[\vv{(\Gamma_k).A_k}/\vv{X_k}]@\id{\Gamma},\Delta_k,y_k:\nu[\vv{(\Gamma_k).A_k}/\vv{X_k}]@\sigma_k\vdash\widehat{A_k}(\vv{t},y_k)[(\xi_k^{\nu[\vv{(\Gamma_k).A_k}/\vv{X_k}]}@\id{\Delta_k}@y_k)/y_k]:A_k[\nu[\vv{(\Gamma_k).A_k}/\vv{X_k}]/X]$}
    \IndE
    \BinaryInfC{$\Gamma,x:\nu[\vv{(\Gamma_k).A_k}/\vv{X_k}]@\id{\Gamma}\vdash\corec^{\nu[\vv{(\Gamma_k).B_k}/\vv{X_k}]}\vv{(\Delta_k,y_k).\widehat{A_k}(\vv{t},y_k)[(\xi_k^{\nu[\vv{(\Gamma_k).A_k}/\vv{X_k}]}@\id{\Delta_k}@y_k)/y_k]}:(\Gamma,x:\nu[\vv{(\Gamma_k).A_k}/\vv{X_k}]@\id{\Gamma})\rat\nu[\vv{(\Gamma_k).B_k}/\vv{X_k}]$}
    \AxiomC{$\ldots$}
    \Inst{$\Gamma,x:\nu[\vv{(\Gamma_k).A_k}/\vv{X_k}]@\id{\Gamma}\vdash\corec^{\nu[\vv{(\Gamma_k).B_k}/\vv{X_k}]}\vv{(\Delta_k,y_k).\widehat{A_k}(\vv{t},y_k)[(\xi_k^{\nu[\vv{(\Gamma_k).A_k}/\vv{X_k}]}@\id{\Delta_k}@y_k)/y_k]}@\id{\Gamma}@x:\nu[\vv{(\Gamma_k).B_k}/\vv{X_k}]$}
  \end{scprooftree}
 \end{landscape}


\bibliographystyle{alpha}
\bibliography{references}
\end{document}
